\documentclass[10pt,shortpaper,twoside,web,logowidth=10.5]{ieeecolor}
\usepackage{generic}
\usepackage{cite}
\usepackage{amsmath,amssymb,amsfonts}
% \usepackage{algorithmic}
\usepackage{graphicx}
\usepackage{textcomp}
\usepackage{stfloats}
\usepackage{graphicx}
\usepackage{booktabs}
\usepackage{float}
\usepackage{amssymb,amsfonts}
\usepackage{listings}
\usepackage{url}
% \usepackage{algorithm,algorithmicx}
\usepackage[ruled,vlined]{algorithm2e}
\usepackage{algpseudocode}
%\usepackage{enumitem}
\usepackage{mathrsfs}
\usepackage{amsmath,bm}
\usepackage[switch]{lineno}
\usepackage[table,xcdraw]{xcolor}
\usepackage{scalerel}
\usepackage{tikz}
\usepackage{multirow}
\usepackage{tipa}

\usepackage{url}
\usepackage{balance}
\usepackage{graphicx}
\usepackage{amsfonts}
\usepackage{fancyhdr}
\usepackage{comment}
\usepackage{times}
\usepackage{amsmath}
\usepackage{changepage}
\usepackage{amssymb}
\usepackage{enumerate}
%\usepackage{algorithm}
% \usepackage{algorithmic}
\usepackage{bm}
\usepackage{subfigure}
\usepackage{calligra}
\usepackage{multirow}
\usepackage{tabularx}
\usepackage{booktabs}
\usepackage{mathtools}
\usepackage{arydshln}
\usepackage{latexsym}
\usepackage{amsmath}
\usepackage{amssymb}
\usepackage{color}
\usepackage{soul}
\usepackage{gensymb}
\usepackage{cite}
\usepackage{color}
\usepackage{cuted}
\usepackage{etoolbox}
\usepackage{capt-of}
\usepackage[flushleft]{threeparttable}
\usepackage{ulem}
\usepackage{float}
\usepackage{booktabs}
\usepackage{bookmark}
\usepackage{hyperref}
% \usepackage[ruled,vlined]{algorithm2e}
\usetikzlibrary{svg.path}

\definecolor{orcidlogocol}{HTML}{A6CE39}
\tikzset{
  orcidlogo/.pic={
    \fill[orcidlogocol] svg{M256,128c0,70.7-57.3,128-128,128C57.3,256,0,198.7,0,128C0,57.3,57.3,0,128,0C198.7,0,256,57.3,256,128z};
    \fill[white] svg{M86.3,186.2H70.9V79.1h15.4v48.4V186.2z}
                 svg{M108.9,79.1h41.6c39.6,0,57,28.3,57,53.6c0,27.5-21.5,53.6-56.8,53.6h-41.8V79.1z M124.3,172.4h24.5c34.9,0,42.9-26.5,42.9-39.7c0-21.5-13.7-39.7-43.7-39.7h-23.7V172.4z}
                 svg{M88.7,56.8c0,5.5-4.5,10.1-10.1,10.1c-5.6,0-10.1-4.6-10.1-10.1c0-5.6,4.5-10.1,10.1-10.1C84.2,46.7,88.7,51.3,88.7,56.8z};
  }
}
\newcommand\orcidicon[1]{\href{https://orcid.org/#1}{\mbox{\scalerel*{
\begin{tikzpicture}[yscale=-1,transform shape]
\pic{orcidlogo};
\end{tikzpicture}
}{|}}}}

\usepackage{hyperref}
\hypersetup{
    colorlinks=true,
    linkcolor=black,
    filecolor=black,      
    urlcolor=black,
    citecolor=black
}

\newcommand{\True}{{\rm{true}}}
\newcommand{\SO}{{\rm{SO}}}
\newcommand{\SL}{{\rm{SL}}}
\newcommand{\est}{{\rm{est}}}
\newcommand{\reff}{{\rm{ref}}}
\newcommand{\tr}{{\rm{tr}}}
\newcommand{\vecc}{{\rm{vec}}}
\newcommand{\dif}{{\rm{d}}}

\def\BibTeX{{\rm B\kern-.05em{\sc i\kern-.025em b}\kern-.08em
    T\kern-.1667em\lower.7ex\hbox{E}\kern-.125emX}}
% \markboth{IEEE/ASME Transactions on Mechatronics}
% \markboth{Submitted to IEEE/ASME Transactions on Mechatronics}
\markboth{}
{J. Jiao \MakeLowercase{\textit{et al.}}: 
LCE-Calib: Automatic LiDAR-Frame/Event Camera Extrinsic Calibration With A Globally Optimal Solution}

\newcommand{\rt}{\textcolor[rgb]{1,0,0}}
\newcommand{\gt}{\textcolor[rgb]{0,1,0}}
\newcommand{\bt}{\textcolor[rgb]{0,0,1}}
\newcommand{\pt}{\textcolor[rgb]{0.466,0.0,0.78}}
\newcommand{\dgt}{\textcolor[rgb]{0.4,0.4,0.4}}
\newcommand{\darkgraytext}{\textcolor[rgb]{0.4,0.4,0.4}}

%%%%%%%%%%%%%%%%%%%%%%%%%%%%%%%%%%%%%%%%%%%%%%
%%%%%%%%%%%%%%%%%%%%%%%%%%%%%%%%%%%%%%%%%%%%%%
%%%%%%%%%%%%%%%%%%%%%%%%%%%%%%%%%%%%%%%%%%%%%%
\begin{document}
\title{\LARGE{LCE-Calib: Automatic LiDAR-Frame/Event Camera Extrinsic Calibration With A Globally Optimal Solution}}

\author{\large{Jianhao Jiao$^{2,4}$, \textit{IEEE Member}, Feiyi Chen$^{2}$, Hexiang Wei$^{2}$, Jin Wu$^{2}$, \textit{IEEE Member}, \\Ming Liu$^{1,2,3}$, \textit{IEEE Senior Member}}
  \thanks{This work was supported by Guangdong Basic and Applied Basic Research Foundation, under project 2021B1515120032, Foshan-HKUST Project no. FSUST20-SHCIRI06C, and Project of Hetao Shenzhen-Hong Kong Science and Technology Innovation Cooperation Zone (HZQB-KCZYB-2020083), awarded to Prof. Ming Liu.}
  \thanks{$^{1}$Robotics and Autonomous Systsms, The Hong Kong University of Science and Technology (Guangzhou), Nansha, Guangzhou, 511400, Guangdong, China.}
  \thanks{$^{2}$Department of Electronic and Computer Engineering, Hong Kong University of Science and Technology, Hong Kong, China. \{jjiao,hweiak,fchenak,jwucp\}@connect.ust.hk, \{eelium\}@ust.hk.}
  \thanks{$^{3}$HKUST Shenzhen-Hong Kong Collaborative Innovation Research Institute, Futian, Shenzhen.}
  \thanks{$^{4}$CWB Institute of Autonomous Driving, Nanshan, Shenzhen.}
  \thanks{\textit{Corresponding author: Ming Liu.}}
}

\maketitle
%\linenumbers



Over the past few years, there has been a significant amount of research focused on studying the ReLU activation function, with the aim of achieving neural network convergence through over-parametrization. However, recent developments in the field of Large Language Models (LLMs) have sparked interest in the use of exponential activation functions, specifically in the attention mechanism.

Mathematically, we define the neural function $F: \R^{d \times m} \times  \mathbb{R}^d \rightarrow \mathbb{R}$ using an exponential activation function. Given a set of data points with labels $\{(x_1, y_1), (x_2, y_2), \dots, (x_n, y_n)\} \subset \mathbb{R}^d \times \mathbb{R}$ where $n$ denotes the number of the data. Here $F(W(t),x)$ can be expressed as $F(W(t),x) := \sum_{r=1}^m a_r \exp(\langle w_r, x \rangle)$, where $m$ represents the number of neurons, and $w_r(t)$ are weights at time $t$. It's standard in literature that $a_r$ are the fixed weights and it's never changed during the training. We initialize the weights $W(0) \in \mathbb{R}^{d \times m}$ with random Gaussian distributions, such that $w_r(0) \sim \mathcal{N}(0, I_d)$ and initialize $a_r$ from random sign distribution for each $r \in [m]$.

Using the gradient descent algorithm, we can find a weight $W(T)$ such that $\| F(W(T), X) - y \|_2 \leq \epsilon$ holds with probability $1-\delta$, where $\epsilon \in (0,0.1)$ and $m = \Omega(n^{2+o(1)}\log(n/\delta))$. To optimize the over-parametrization bound $m$, we employ several tight analysis techniques from previous studies [Song and Yang arXiv 2019, Munteanu, Omlor, Song and Woodruff ICML 2022]. 

 


\begin{IEEEkeywords}
  \small{Machine and Computer Vision, Unmanned Autonomous Systems}
\end{IEEEkeywords}


\section{Introduction}
\label{sec:introduction}
% \begin{itemize}
%     % Diffusion of FL
%     \item {\st{Diffusion of FL}}
%     % Security threats to FL
%     \item {\st{Security threats to FL with particular focus on model poisoning}}
%     % Limitations of existing countermeasures
%     \item {\st{Current countermeasures (e.g., KRUM) and their limitations}}
%     % Proposed method and its advantages
%     \item {\st{Intuitive description of the proposed method and its difference (i.e., advantages) w.r.t. state of the art}}
%     % Main contributions
%     \item {\st{Summary of the main contributions of this work}}
%     % Paper's structure and organization
%     \item {\st{Paper's structure and organization}}
% \end{itemize}

% Diffusion of FL
Recently, {\em federated learning} (FL) has emerged as the leading paradigm for training distributed, large-scale, and privacy-preserving machine learning (ML) systems~\cite{mcmahan2017googleai,mcmahan2017aistats}. 
The core idea of FL is to allow multiple edge clients to collaboratively train a shared, global model without disclosing their local private training data.
%Specifically, an FL system consists of a central server and many edge clients; 
A typical FL round involves the following steps: {\em(i)} the server randomly picks some clients and sends them the current, global model; {\em(ii)} each selected client locally trains its model with its own private data; then, it sends the resulting local model to the server;\footnote{Whenever we refer to global/local model, we mean global/local model {\em parameters}.} {\em(iii)} the server updates the global model by computing an \emph{aggregation function}, usually the average (FedAvg), on the local models received from clients.
% \begin{enumerate}
%     \item[{\em(i)}] the server sends the current, global model to the clients and appoints some of them for training;
%     \item[{\em(ii)}] each selected client locally trains its copy of the global model with its own private data; then, it sends the resulting local model back to the server;\footnote{Whenever we refer to global/local model, we mean global/local model {\em parameters}.}
%     \item[{\em(iii)}] the server updates the global model by computing an \emph{aggregation function} on the local models received from clients (by default, the average, also referred to as FedAvg~\cite{mcmahan2017aistats}).
% \end{enumerate}
This process goes on until the global model converges. %(e.g., after a certain number of rounds or other similar stopping criteria).
%\\
% The advantages of FL over the traditional, centralized learning paradigm are undoubtedly clear in terms of flexibility/scalability (clients can join/disconnect from the FL network dynamically), network communications (only model weights\footnote{We will use \textit{parameters} and \textit{weights} interchangeably.} are exchanged between clients and server), and privacy (each client's private training data is kept local at the client's end and not uploaded to the server).
\\
% Security threats to FL
%However, the growing adoption of FL also raises security concerns~\cite{costa2022covert}, particularly about its confidentiality, integrity, and availability.
Although its advantages over standard ML, FL also raises security concerns~\cite{costa2022covert}. %, particularly about its confidentiality, integrity, and availability~\cite{costa2022covert}.
% OLD, LONG VERSION
% Indeed, some work deals with privacy leakage that may expose the local data of some clients~\cite{melis2019sp}. 
% A large body of work, instead, investigates attacks that usually aim to detriment the predictive accuracy of the learned global model. For instance, \emph{data poisoning} attacks achieve this goal by letting an adversary pollute the training set of some corrupt FL clients with maliciously crafted examples~\cite{jagielski2018sp}.
% Similarly, in \emph{model poisoning} the attacker attempts to tweak the global model weights~\cite{bhagoji2019pmlr} by directly perturbing the local model's weights of some infected FL clients before these are sent to the central server for aggregation, usually via so-called Byzantine attacks. 
% It turns out that Byzantine model poisoning attacks severely impact standard FedAvg; therefore, more robust aggregation functions must be designed to make FL systems secure.
Here, we focus on \emph{untargeted model poisoning} attacks~\cite{bhagoji2019pmlr}, where an adversary attempts to tweak the global model weights %\footnote{We will use the terms \textit{parameters} and \textit{weights} interchangeably.} 
by directly perturbing the local model's parameters of some infected clients before these are sent to the central server for aggregation.
In doing so, the adversary aims to jeopardize the global model \textit{indiscriminately} at inference time.
Such model poisoning attacks severely impact standard FedAvg; therefore, more robust aggregation functions must be designed to secure FL systems.
\\
% In this paper, we focus on designing a novel robust aggregation scheme at the server's end to contrast the effect of Byzantine model poisoning attacks.
%
% Current countermeasures and their limitations
%Several countermeasures have been proposed in the literature to combat model poisoning attacks on FL systems.
% Some methods use simple statistics more robust than plain average to smooth the impact of malicious updates (e.g., Trimmed Mean and FedMedian~\cite{yin2018icml}). 
% Other defenses implement outlier detection techniques to discard malicious updates from the aggregation performed at the server's end. Those are either based on heuristics (e.g., Krum/Multi-Krum~\cite{blanchard2017nips} and Bulyan~\cite{mhamdi2018pmlr}) or data-driven approaches (e.g., K-means clustering~\cite{shen2016acm} or DnC via spectral analysis~\cite{shejwalkar2021ndss}). 
% Finally, some strategies rely on a centralized ``source of trust'' to spot potential malicious updates (e.g., FLTrust~\cite{cao2020fltrust}).
% Several countermeasures have been proposed in the literature to combat model poisoning attacks on FL systems, i.e., to discard possible malicious local updates from the aggregation performed at the server's end. 
% These techniques range from simple statistics more robust than plain average (e.g., Trimmed Mean and FedMedian~\cite{yin2018icml}) to outlier detection heuristics (e.g., Krum/Multi-Krum~\cite{blanchard2017nips} and Bulyan~\cite{mhamdi2018pmlr}) or data-driven approaches (e.g., spectral analysis via K-means clustering~\cite{shen2016acm} or spectral analysis), or methods based on ``source of trust'' (e.g., FLTrust~\cite{cao2020fltrust}).
% OLD, LONG VERSION
%Several countermeasures have been proposed in the literature to combat Byzantine model poisoning attacks on FL systems.
% Descriptive statistics
% For example, Trimmed Mean and FedMedian aggregate local model updates using more robust statistics than standard average~\cite{yin2018icml}.
%
% % Heuristics for outlier detection
% Many existing Byzantine-resilient strategies implement some outlier detection heuristics to discard the model updates sent by potentially malicious clients from the input of the aggregation function.
% One of the most popular heuristics is Krum~\cite{blanchard2017nips}.
% This strategy tries to mitigate the impact of Byzantine attacks by selecting as a global model the local model with the smallest sum of Euclidean distances to {\em all} the other local models.
% Although powerful, Krum requires the server to know (or, at least, estimate) the number of malicious FL clients upfront, which is generally impossible in a realistic attack scenario. %
% Moreover, Krum may become ineffective for complex, high-dimensional model parameter spaces due to the curse of dimensionality.
% Bulyan~\cite{mhamdi2018pmlr} tries to overcome this issue by combining Krum with a variant of Trimmed Mean.
% % Data-driven outlier detection
% Other strategies use data-driven outlier detection techniques -- e.g., via K-means clustering~\cite{shen2016acm} -- to spot potential malicious local model updates. 
% %For instance, Shen et al. propose to cluster local model updates with K-means and thus identify outliers.
%
% % Other techniques
% As far as the server is concerned, any local model received can be from a potential malicious client. 
% FLTrust~\cite{cao2020fltrust} assumes the server acts as a client, i.e., trains a local model on an additional {\em trustworthy} dataset at the server's end and compares it against all the local models from other clients. 
% This way, the server can rely on some ``source of trust'' when discarding potentially malicious clients.
%\\
% Limitations of existing Byzantine-resilient strategies
Unfortunately, existing defense mechanisms either rely on simple heuristics (e.g., Trimmed Mean and FedMedian by~\cite{yin2018icml}) or need strong and unrealistic assumptions to work effectively (e.g., foreknowledge or estimation of the number of malicious clients in the FL system, as for Krum/Multi-Krum~\cite{blanchard2017nips} and Bulyan~\cite{mhamdi2018pmlr}, which, however, cannot exceed a fixed threshold).
Furthermore, outlier detection methods using K-means clustering~\cite{shen2016acm} or spectral analysis like DnC~\cite{shejwalkar2021ndss} do not directly consider the temporal evolution of local model updates received.
Finally, strategies like FLTrust~\cite{cao2020fltrust} require the server to collect its own dataset and act as a proper client, thereby altering the standard FL protocol.
\\
% OLD, LONG VERSION
% Overall, existing Byzantine-resilient strategies are either simple heuristics (e.g., FedMedian) or, if they are more complex, they rely on strong and unrealistic assumptions to work effectively (e.g., knowing the number of malicious clients in the FL system in advance, as for Krum and alike).
% Furthermore, data-driven outlier detection methods do not consider the temporary evolution of local model updates received (e.g., K-means clustering). 
% Finally, strategies like FLTrust requires the server to collect its own dataset and act as a proper client, thereby altering the standard FL protocol.
%
% Description of the proposed method
This work introduces a novel pre-aggregation \textit{filter} robust to untargeted model poisoning attacks. Notably, this filter $(i)$ operates without requiring prior knowledge or constraints on the number of malicious clients and $(ii)$ inherently integrates temporal dependencies. 
The FL server can employ this filter as a preprocessing step before applying \textit{any} aggregation function, be it standard like FedAvg or robust like Krum or Bulyan.
Specifically, we formulate the problem of identifying corrupted updates as a multidimensional (i.e., matrix-valued) time series anomaly detection task. 
The key idea is that legitimate local updates, resulting from well-calibrated iterative procedures like stochastic gradient descent (SGD) with an appropriate learning rate, show \textit{higher predictability} compared to malicious updates. This hypothesis stems from the fact that the sequence of gradients (thus, model parameters) observed during legitimate training exhibit regular patterns, as validated in Section~\ref{subsec:intuition}. %until convergence. 
%This regularity may be more pronounced for smooth convex loss functions, but it can still be captured within an appropriate time window, even for more complex and convoluted loss surfaces. 
%We provide evidence of this claim in Appendix~B, where we show that the average mutual information (i.e., ``predictability''), calculated over pairs of legitimate model updates sent at different FL rounds, is significantly higher than the corresponding computation for a malicious client.
\\
Inspired by the matrix autoregressive (MAR) framework for multidimensional time series forecasting~\cite{chen2021je}, we propose the FLANDERS ({\em \textbf{F}ederated \textbf{L}earning meets \textbf{AN}omaly \textbf{DE}tection for a \textbf{R}obust and \textbf{S}ecure}) filter.
The main advantages of FLANDERS over existing strategies like FLDetector~\cite{zhao2020multivariate} are its resilience to large-scale attacks, where $50\%$ or more FL participants are hostile, and the capability of working under realistic non-iid scenarios.
We attribute such a capability to two key factors: $(i)$ FLANDERS works without knowing a priori the ratio of corrupted clients, and $(ii)$ it embodies temporal dependencies between intra- and inter-client updates, quickly recognizing local model drifts caused by evil players. Below, we summarize our main contributions:

\begin{itemize}
\item[{\em(i)}]
We provide empirical evidence that the sequence of models sent by legitimate clients is more predictable than those of malicious participants performing untargeted model poisoning attacks.
\\
\item[{\em(ii)}] 
We introduce FLANDERS, the first pre-aggregation filter for FL robust to untargeted model poisoning based on multidimensional time series anomaly detection.
\\
\item[{\em(iii)}] 
We integrate FLANDERS into Flower,\footnote{\scriptsize{\url{https://flower.dev/}}} a popular FL simulation framework for reproducibility.
\\
\item[{\em(iv)}] 
We show that FLANDERS improves the robustness of the existing aggregation methods under multiple settings: different datasets, client's data distribution (non-iid), models, and attack scenarios.
\\
\item[{\em(v)}] 
We publicly release all the implementation code of FLANDERS along with our experiments.\footnote{\scriptsize{\url{https://anonymous.4open.science/r/flanders_exp-7EEB}}}
\end{itemize}

% Paper's structure and organization
The remainder of the paper is structured as follows. %some related work and the current state-of-the-art solutions to security issues that FL entails. 
Section~\ref{sec:background} covers background and preliminaries. 
In Section~\ref{sec:related}, we discuss related work.
Section~\ref{sec:problem} and Section~\ref{sec:method} describe the problem formulation and the method proposed. % to tackle it. 
Section~\ref{sec:experiments} gathers experimental results. %, and Section~\ref{sec:limitations} discusses some limitations of this work.
Finally, we conclude in Section~\ref{sec:conclusion}.
 %discusses the limitations of this work and draws future research directions.
%reports conclusions and draws perspectives for future research directions.

%%%%%%% OLD %%%%%%%
%to overcome the resilience of Byzantine failures in distributed Stochastic Gradient Descent computations. 
% The strength of Krum is its time complexity, which is linear in the gradient dimension. 
% However, the robustness of the approach is guaranteed for gradient-based learning applications only when the majority of the clients are not compromised. 
% Besides, the aggregation mechanism of Krum, as well as that of similar methods, is robust from a coarse-grained perspective and does not provide solutions to errors and perturbations that may occur at inference time.
%A related approach to~\cite{blanchard2017nips} is the work of Su et al.~\cite{su2016dc}. Here, the authors propose an iterated approximate agreement to tackle a multi-layer scenario attacked by Byzantine agents. 
%However, the method works efficiently on the sole discrete context and it is inapplicable to continuous state environments.
%\gabri{Maybe, we should just talk about the main limitations of existing countermeasures without digging into their details (or, we can just mention Krum as this is the most popular one). I will move the description of all these methods to the Related Work section.}
\section{Related work}
% There is extensive recent work on speeding up analytical queries due to the need for consistent execution times in the face of the explosive growth in the volume of available data.
% In this section, we divide existing work into two categories: maintaining data freshness in MVs (\Cref{sec:server_side}) and optimizations for minimizing ad-hoc query latency (\Cref{sec:client_side}).

% \subsection{Maintaining Data Freshness in MVs}
% \label{sec:server_side}
% There exists a variety of data warehousing applications aimed at supporting low-latency analytical queries on fresh data.
% In particular, these applications require efficiency in the propagation of newly ingested data into downstream MVs.
 
\mypara{Efficient MV Refresh}
Incremental view maintenance (IVM) aims to update MVs to reflect newly ingested data, taking advantage of already computed results to perform the update in a manner more efficient than computing from scratch (full refresh)
~\cite{ahmad2012dbtoaster,mcsherry2013differential,armbrust2013generalized,zeng2016iolap, palpanas2002incremental, griffin1995incremental, agiwal2021napa, braun2015analytics}. 
There is an abundance of work in IVM, including incremental updates on duplicate values~\cite{griffin1995incremental}, non-distributive aggregate functions~\cite{palpanas2002incremental}, higher-order views~\cite{ahmad2012dbtoaster}, and sliding windows~\cite{braun2015analytics}. 
More recent works also investigate the scalability aspect of IVM, proposing scale-independent updates~\cite{armbrust2013generalized} and sampled views~\cite{zeng2016iolap}. Since \system is applicable to arbitrary SQL statements, \system is orthogonal to and is fully compatible with existing IVM techniques.

\mypara{MV Refresh Scheduling}
There exist works on scheduling the refresh of a MV set focusing on resolving cyclic dependencies~\cite{folkert2005optimizing}, minimizing weighted average staleness~\cite{golab2009scheduling}, and prioritizing MVs with the highest speedups on predicted future queries~\cite{ahmed2020automated}.
\system's scheduling to speed up the end-to-end refresh of the MV set is not addressed in existing works.

\mypara{DAG Workflow Scheduling}
The execution of workloads consisting of individual jobs with acyclic dependencies is a well-studied topic~\cite{apacheoozie,sparkdag,marchal2018parallel,bathie2020revisiting,baruah2022ilp}; many of these techniques can be applied to MV refresh runs studied in this paper.
Existing workflow scheduling systems such as Apache Oozie~\cite{apacheoozie}, Apache Airflow~\cite{airflow}, and Spark DAG scheduler~\cite{sparkdag} automate the execution of user-defined workflows following a topological order.
There are recent works aimed at finding more optimal execution orders in terms of peak memory usage~\cite{marchal2018parallel, bathie2020revisiting} and execution time on parallel platforms~\cite{baruah2022ilp}.
While \system is designed for use with MV refresh runs/workloads, our technique on joint scheduling and optimization can be reasonably applied to general workloads as a possible future direction.

% \paragraph{Incremental MV indexing}
% Update-optimized indices such as the log-structured merge-trees (LSM)~\cite{o1996log} are used for indexing MVs due to frequent updates induced by data ingestion~\cite{gupta2016mesa,agiwal2021napa}.
% \system is orthogonal to indexing: \system is capable of efficiently performing MV refresh runs regardless of whether the individual MVs are indexed or not.

% \subsection{Ad-hoc Query Latency Reduction}
% \label{sec:client_side}

% The minimization of ad-hoc analytical query response times is a well-studied topic due to latency being negatively correlated with the productivity of a data analyst during a data analysis session~\cite{liu2014effects}.
% Sessions are commonly conducted within visualization systems that contain a variety of optimization techniques to ensure that query response times fall within a certain latency tolerance.

% \mypara{Data prefetching}
% Data is often loaded into memory on a by-need basis in visualization systems to minimize interference with user-issued query computations~\cite{mani2017effective,xin2021enhancing,galakatos2017revisiting, yan2020auto, battle2016dynamic, crotty2016case, jalaparti2018netco}. 
% Query-time data retrieval can be significantly expedited by anticipating the data usage of the user in future queries and pre-loading the data into memory during the downtime between user queries (`think time'). SMART~\cite{mani2017effective} prefetches data for modified versions of current user-issued queries with different filters and dimensions. A-WARE~\cite{crotty2016case} maintains a sample store constantly refined through ingesting data based on speculations of future plots.
% ForeCache~\cite{battle2016dynamic} uses an SVM to predict the user's current analysis phase and accordingly prefetches data tiles partitioned based on different numerical values. NetCo predicts future queries via log analysis, and solves an ILP formulation to prefetch data to maximize the number of SLO-meeting queries~\cite{jalaparti2018netco}.
% In the case of MV refresh workloads, `think time' is nonexistent as individual MVs are refreshed back-to-back, rendering data prefetching techniques non-applicable.

\mypara{Intermediate Data Caching}
Some existing data visualization systems cache user-defined variables to support the typical incremental construction of data visualizations~\cite{zgraggen2016progressive, eichmann2020idebench} during data analysis sessions~\cite{jupyter, rstudio, colab}. 
Recent work proposes a management scheme for these cached variables under a memory constraint that greedily keeps variables with the highest estimated time savings based on predicted future user behavior via neural networks~\cite{xin2021enhancing}.
While useful for data visualization, a greedy approach to memory management fails to achieve satisfactory results compared to \system.

\mypara{Intermediate Result Reuse}

There exist works on storing intermediate results from computations to speedup future computations~\cite{yang2018intermediate, dursun2017revisiting, nagel2013recycling, michiardi2019memory, galakatos2017revisiting}.
Studied topics include the identification of reuse opportunities by finding overlaps in computation graphs of successive jobs~\cite{yang2018intermediate, michiardi2019memory},
selective storage under a space constraint with heuristics such as reuse probability~\cite{dursun2017revisiting}, expected savings~\cite{yang2018intermediate}, and recompute-storage cost difference~\cite{nagel2013recycling},
and rewriting incoming jobs to take advantage of stored intermediates~\cite{galakatos2017revisiting}.
These works share similarity with \system in their selection of items to store under a memory constraint, however, \system's problem setting requires it to uniquely consider the joint (re)ordering of job executions along with the selection of items.

% work that considers both job execution (re)order as well as intermediate result caching with a bounded amount of memory. but notably lack the joint aspect of \system and cannot be used to achieve immediate speedup on an incoming MV refresh run if no intermediates are stored beforehand. 

\mypara{Incremental Query Processing} Incremental processing (IQP) is useful for cases where not all data required for a query is immediately available. Similar to online aggregation~\cite{hellerstein1997online}, initial results of a query are computed on a subset of required data and progressively refined as the rest of the required data arrives in a predictable pattern~\cite{tang2019intermittent,wangtempura}. Tang et al. propose a dynamic programming formulation to pick intermediate states to store in memory given a limited memory budget~\cite{tang2019intermittent}. Tempura rewrites the query plan for more efficient execution based on predicted data arrival patterns~\cite{wangtempura}. While similarities exist between the problem setting of IQP and \system, such as management of bounded memory, \system notably includes additional joint optimization for the order of MV updates.

% \paragraph{Sampling}
% Sampling has seen wide use in visualization systems for reducing the computation time of ad-hoc queries by computing an approximate result over a subset of data as exact results are not always required by the user~\cite{crotty2016case, mani2017effective, zgraggen2014panoramicdata, kraska2021northstar, galakatos2017revisiting, kandula2016quickr}. 
% Commonly studied topics in sampling for ad-hoc queries include complex query sampling~\cite{kandula2016quickr}, rare event aggregation~\cite{kraska2021northstar, galakatos2017revisiting}, and maintaining consistency between related sampled visualizations~\cite{zgraggen2014panoramicdata}.
% Sampling server-side at the MV level compromises the assumptions of downstream applications and is thus not considered in \system.

% \paragraph{Progressive visualization}
% The latency tolerance for time-consuming queries can be circumvented by presenting a partially-computed visualization to the user within the tolerance, which is then incrementally refined until it is fully accurate~\cite{rahman2017ve, zgraggen2016progressive, crotty2015vizdom, kraska2021northstar, kamat2017infiniviz}.
% Example plots which benefit from progressive visualization include bar charts~\cite{kamat2017infiniviz} and heatmaps~\cite{rahman2017ve}.
% Similar to sampling, study on this topic is orthogonal to \system as pushing out partially-updated MVs compromises downstream assumptions.
\section{Pipeline inspection exemplar}
\label{sec:problem_statement}
This section describes the use case and system architecture, the two
system layers are detailed in Sections~\ref{sec:managed} and
\ref{sec:managing_subsystem}.

\subsection{Use case description}
\label{sec:use_case_description}
The use case in this exemplar is about an AUV inspecting pipelines located on a seabed.
Its mission consists of two sequential tasks, $(T1)$~searching for the pipeline, then $(T2)$~simultaneously following and inspecting the pipeline. 


When performing its mission, the AUV is subject to two sources of
uncertainty that could trigger self-adaptation: $(U1)$~thruster
failures and $(U2)$~changes in water visibility. $U1$~arises from the
possibility of the AUV's thrusters failing at runtime, which may cause
the AUV to move unexpectedly. This is relevant for both $T1$ and $T2$. To overcome $U1$, the managed subsystem of the AUV
contains functional alternatives. When
one or more thrusters fail, it is possible to enter a recovery state
in which the thrusters are recovered. 
$U2$
influences the maximum distance at which the AUV can visually perceive
objects. This is relevant for $T1$, higher water visibility
allows the AUV to search for the pipeline at higher altitudes,
resulting in a larger field of view and the possibility of discovering the pipeline faster. On the other hand, if the water
visibility is low, the AUV has to move closer to the seabed to
search for the pipeline, which limits its field of view and therefore
may lead to a longer time to discover the pipeline. Thus, changing the
altitude of the AUV provides functional alternatives for dealing with $U2$.

This exemplar focuses on the problem of overcoming $U1$ and $U2$ using
a self-adaptation logic, implemented by a managing subsystem, that
can be extended and reused for other sources of uncertainty.
The managing subsystem shall overcome $U1$ by recovering the failed
thrusters at runtime, and $U2$ by adapting the maximum altitude for
the path generator algorithm according to the measured water
visibility.  Thus, by reacting to $U1$ and $U2$, the managing
subsystem increases the reliability and performance of the system.

For the feasibility of the exemplar, the use case was simplified while still allowing for a worthwhile application of self-adaptation to an AUV.  
 It is important to highlight that a
realistic operation of an AUV used for pipeline inspection would include
steps that are related to pre-dive, launching and recovery, human
interaction, and intermediary missions that are necessary to enable
the inspection.  Furthermore, there are several sources of uncertainty
not considered here, including ocean dynamics, sensor failures, and
battery duration.

\subsection{System Architecture}

To accomplish the mission described in Section~\ref{sec:use_case_description}, the managed subsystem requires the
functions represented in Fig.~\ref{fig:func_arch}. $T1$ requires the
functions \kw{Control Motion}, \kw{Maintain Motion},
\kw{Localization}, \kw{Detect Pipeline}, \kw{Generate Search
  Path}, and \kw{Coordinate Mission}, while $T2$ requires the functions \kw{Control Motion},
\kw{Maintain Motion}, \kw{Localization}, \kw{Detect Pipeline},
\kw{Follow Pipeline}, \kw{Inspect Pipeline}, and \kw{Coordinate Mission}. During runtime, the
functions must be activated and deactivated according to the task
being performed.


\begin{figure}[h]
    \centering
    \includegraphics[width=0.8\linewidth]{figures/FunctionalHierarchy.pdf}
    \caption{Managed Subsystem's Functional Hierarchy}
    \label{fig:func_arch}
\end{figure}
%

To overcome the uncertainties $U1$ and $U2$, a managing subsystem
requires the functionalities to \emph{monitor} the environment and the
managed subsystem's internal state, \emph{reason} about it, and
\emph{execute} the managed subsystem's reconfiguration.


\begin{figure}
    \centering
    \includegraphics[width=\linewidth]{figures/GeneralArchitecture_small.pdf}
    \caption{System Architecture}
    \label{fig:system_architecture}
\end{figure}

The required functions of the managed and managing subsystems
are realized as depicted in Fig.~\ref{fig:system_architecture}. The
managed subsystem is detailed in Section~\ref{sec:managed} and the
managing subsystem in Section~\ref{sec:managing_subsystem}.
It is important to mention that managed subsystem functions
\kw{Control Motion} and \kw{Localization} are achieved by \kw{ArduSub}, and the function \kw{Inspect Pipeline} is not realized since
the actual inspection of the pipeline is not the focus of this
work. It is also important to highlight that this exemplar implements
the function to \emph{reason} about the managing subsystem with Metacontrol to provide a
baseline for future research. However, it can be replaced with other
solutions, as long as they are compatible with the monitor and execute
interfaces, as described in Section~\ref{sec:managing_requirements}.

\section{Measurement Processing}
\label{sec:measure_process}

\begin{figure}[t]
  \centering
  \subfigure[]
  {\label{fig:mp_pipeline_result_cloud}\centering\includegraphics[width=0.19\textwidth]{figure/methodology/mp_cloud_3-crop}}
  \subfigure[]
  {\label{fig:mp_pipeline_result_frame_image}\centering\includegraphics[width=0.14\textwidth]{figure/methodology/mp_frame_image-crop}}
  \subfigure[]
  {\label{fig:mp_pipeline_result_event_image}\centering\includegraphics[width=0.14\textwidth]{figure/methodology/mp_event_image-crop}}
  \caption{Features of the checkerboard are extracted from (a) the point cloud, (b) frame image, and (c) reconstructed image from events.}
  \label{fig:mp_pipeline_result}
\end{figure}

This section explains how raw sensor data from LiDARs, frame cameras, and event cameras are preprocessed before estimating their extrinsics.
We can take advantage of the structural prior of the checkerboard for reliable feature selection and matching.
Fig. \ref{fig:mp_pipeline_result} shows the feature extraction results.

\subsection{Automatic Feature Extraction From LiDARs}
\label{sec:methodology_fsl}
We are interested in extracting three types of features from the checkerboard point cloud: edge points, planar points, and planar coefficients.
These features are useful in LiDAR-camera data association.
Since we work on mechanical scanning LiDARs, the following sections use \textit{ring} to denote the points set from the same emitter.
We propose a two-stage method for automatic feature extraction.

%%%%%%%%%%%%%%%%%%%%%%%%%%%%%%%%%%%%%%%%%%%%
\subsubsection{Checkerboard Point Selection}
\label{sec:methodology_fsl_selection}
We need to select board points from the raw point cloud.
We have these observations to design the point selection method:
\textit{1)} the width and height of the board are known as a prior;
\textit{2)} points from the same ring form a straight line segment if they lie on the board;
\textit{3)} points from different rings form a planar patch if they lie on the board and stay near.

Firstly, we compute the pitch angle $\phi = \arctan(\frac{z}{\sqrt{x^2 + y^2}})$ for each point and split out points into different rings.
Some LiDARs directly provide the ``ring'' as a property of each point.
We then cluster each ring into multiple line segments based on the distance of two consecutive points.
The principle component of each line segment is computed by PCA: $[\sigma_1, \sigma_2, \sigma_3]$ where $\sigma_1 \geq \sigma_2 \geq \sigma_3$. The curvature is measured as $\sigma_1/\sigma_2$.
We keep line segments if their curvature is bigger than the predefined threshold $\mu_1$ and the length of the line segment smaller than an empirical threshold $\mu_2$.
Secondly, we utilize the DBScan algorithm \cite{kriegel2011density} to segment the remaining points as several clusters. Small clusters are then removed.
Thirdly, we verify each cluster by registering it with a board point template.
We compute the registration error for each cluster and keep the only one with the lowest error. The selected cluster is regarded as the candidate board points.

%%%%%%%%%%%%%%%%%%%%%%%%%%%%%%%%%%%%%%%%%%%%
\subsubsection{Checkerboard Feature Extraction}
\label{sec:methodology_fsl_feature_extraction}

After obtaining board points, we use the RANSAC-based plane fitting method \cite{rusu20113d} to compute the normal vector $\bm{n}^{l}$, filter out outliers, and obtain a set of planar points $\mathcal{P}^{l}_{pl}$ (inliers).
Due to uncorrected bias and noise of LiDARs' measurements, we observe that board points form a $3cm$-thick plane approximately.
To reduce this effect in the subsequent calibration process, we project each planar point $\bm{p}$ onto the fitted plane as
\begin{equation}
  \label{equ:projection_on_plane}
  \bm{p}_{proj}^{l}
  =
  \bm{p}^{l} - a\bm{n}^{l},
  \ \ \
  a = \bm{n}^{l}\cdot(\bm{p}^{l} - \bm{g}^{l}),
\end{equation}
where $\bm{g}^{l}$ is a point on the plane.
The edge points $\mathcal{P}^{l}_{edge}$ are then extracted as the starting and ending points of each ring.

% We will show how this projection can improve the accuracy in Section \ref{sec:experiment}.}

% pts_on_plane = lidar_board_edge_pts(1:3, 1);
% pts_on_plane(3) = -((d + n(1:2)' * pts_on_plane(1:2)) / n(3));

% for ptid = 1 : size(lidar_board_pts, 2)
%   pt = lidar_board_pts(1:3, ptid);
%   v = pt - pts_on_plane;
%   pt_onboard = pts_on_plane + v - (v' * n * n);
%   lidar_board_pts(1:3, ptid) = pt_onboard;
% end

% In contrast, the determination of corners $\mathcal{P}^{l}_{cor}$ is a little complex. 
% A primary solution was proposed in \cite{zhou2018automatic}, where boundary equations are found via the RANSAC-based line fitting algorithm.
% But this scheme presents two limitations: 
% 1) the RANSAC-based fitting process only succeeds if more than three boundary points are provided, 
% and 2) the fitted lines are slightly unreliable since boundary points suffer from large non-Gaussian noise \cite{pomerleau2012noise}.

% Instead, our method fully utilizes the prior knowledge of the checkerboard.
% We formulate the point-to-edge distance minimization problem by computing the optimal transformation from the template board to data $\bm{x}^{l}_{b}$.
% The coarse initialization of orientation and translation of the board is insufficient for a nonlinear optimization problem. 
% We thus employ the global search method based on the scatter-search mechanism \cite{ugray2007scatter} to generate multiple start points around this initial value. 
% For each start point, the corresponding line of each boundary point can be found by the Nearest Neighbor Search (NNS), and we then minimize the sum of all point-to-edge residuals to optimize the transformation as iteratively
% \begin{equation}
% \begin{aligned}
%   \label{equ:methodology_board_pose}
%   &\hat{\bm{x}}^{l}_{b}
%   =
%   \underset{\bm{x}^{l}_{b}}{\arg\min}
%   \sum_{\bm{p}\in{\mathcal{P}_{edge}^{l}}}{}
%   \rho
%   \big(
%     ||\bm{r}_{edge}(\bm{x}^{l}_{b}, \bm{p}, L)||^{2}_{\bm{\Sigma}_{\bm{p}}}
%   \big).
% \end{aligned}
% \end{equation}

% Finally, we can select the optimal result with the lowest error from a set of candidate transformations. It helps us to avoid local minima caused by incorrect data association.
% We can use this resulting transformation to transform original corner points to obtain the observed board corners $\mathcal{P}^{l}_{cor}$.

%%%%%%%%%%%%%%%%%%%%%%%%%%%%%%%%%%%%%%%%%%%%
\subsubsection{Checkerboard Tracking}

Estimating the board's position allows us to track the board at the next frame.
The time-consuming point selection stage (costs $50-60ms$) is skipped.
We create a virtual bounding box around the checkerboard with a slightly larger size.
In the next frame, only points inside the box are kept.
This tracking stage also improves the success rate of the board point selection.

%%%%%%%%%%%%%%%%%%%%%%%%%%%%%%%%%%%%%%%%%%%%
\subsection{Automatic Feature Extraction From Frame Cameras}
\label{sec:methodology_fsc}

The coefficients of the board observed by the frame camera, denoted by $[\bm{n}^{c},d^{c}]^{\top}$, can be easily obtained.
We use the off-the-shelf software to detect the inner patterns of the checkerboard automatically.
The board's pose $SE_{3}(\bm{R}^{c}_{b}, \bm{t}^{c}_{b})$ is estimated by the \textit{QPEP-PnP} algorithm, where we minimize the point-wise distance between 3D inner patterns and their corresponding 2D points.
Along the transformed boundaries of the board, we can generate many ``fake'' edge points $\mathcal{P}_{edge}^{c}$.
We can also obtain the covariance matrix $\bm{\Sigma}_{\bm{\delta\phi}}$ of the checkerboard's rotation from the \textit{QPEP-PnP} algorithm.
Based on the uncertainty representation of 6-DoF poses in Section \ref{sec:ps_notation}, we can propagate the uncertainty of the rotated normal vector $\bm{n}^{c}\approx\bar{\bm{R}}^{c}_{b}(\bm{I} + \delta\phi^{\wedge})\bm{n}^{b}$ as
\begin{equation}
  \begin{split}
    % \bm{n}^{c}
    % &\approx
    % \bar{\bm{R}}^{c}_{b}(\bm{I} + \delta\phi^{\wedge})\bm{n}^{b}, \\
    \bm{\Sigma}_{\bm{n}^{c}}
    \approx
    [\bar{\bm{R}}^{c}_{b}(\bm{n}^{b})^{\wedge}]
    \bm{\Sigma}_{\bm{\delta\phi}}
    [\bar{\bm{R}}^{c}_{b}(\bm{n}^{b})^{\wedge}]^{\top},
  \end{split}
\end{equation}
where the derived covariance will be used in the optimization objective in Section \ref{sec:ext_calibration}.
% The derivation is detailed in the supplementary material \cite{}.
% \begin{equation}
%   % \underset{\bm{R}^{c}_{b}\in SO(3),\ \bm{t}^{c}_{b}\in\mathbb{R}^{3}}{\arg\min}
%   % \sum_{\bm{p}\in\mathcal{P}_{pat}^{b}}
%   % ||\pi(\bm{p}, \bm{R}^{c}_{b}, \bm{t}^{c}_{b}) - \bm{u}||^{2}_{\bm{\Sigma}},
%   \underset{\bm{R}^{c}_{b}\in SO(3),\bm{t}^{c}_{b}\in\mathbb{R}^{3}}{\arg\min}
%   \sum_{\bm{p}\in\mathcal{P}_{pat}^{b}}
%   f_{pnp}(\bm{p}, \bm{R}^{c}_{b}, \bm{t}^{c}_{b}).
%   % ||\bm{r}(\bm{p}, \bm{R}^{c}_{b}, \bm{t}^{c}_{b}) - \bm{u}||^{2}_{\bm{\Sigma}},  
% \end{equation}
% where $\bm{u}$ is the corresponding 2D projected point of $\bm{p}$ and $\pi(\cdot)$ is the camera projection function.}


% The transformation and planar parameters of the checkerboard are calculated by solving the Perspective-n-Point (PnP) problem \cite{Hartley2004}. 
% Thus, we can determine corner features in the camera coordinate system as $\mathcal{P}^{c}_{cor}$. 


%%%%%%%%%%%%%%%%%%%%%%%%%%%%%%%%%%%%%%%%%%%%
\subsection{Image Reconstruction From Events}
\label{sec:methodology_ire}
\subsubsection{Event Data}
An event camera has independent pixels that respond to logarithmic intensity change $L$.
In a noise-free scenario, an event $\bm{e}_{k}=[\bm{u}_{k}^{\top},t_{k},p_{k}]^{\top}$ is triggered at pixel $\bm{u}_{k} = [x_{k},y_{k}]^{\top}$ and time $t_{k}$ as soon as the logarithmic intensity increment reaches a contrast threshold $\pm C$ since the last event at the pixel, i.e.
\begin{equation}
  \Delta L(\bm{u}_{k},t_{k})
  \doteq
  L(\bm{u}_{k}, t_{k}) - L(\bm{u}_{k}, t_{k}-\Delta t_{k})
  \geqslant
  p_{k}C,
\end{equation}
where $C > 0$, $\Delta t_{k}$ is the time elapsed since the last event at the same pixel, and the polarity $p_{k}\in\{+1,-1\}$ is the sign of the intensity change \cite{gallego2020event}.


%%%%%%%%%%%%%%%%%%%%%%%%%%%%%%%%%%%%%%%%%%%%
\begin{algorithm}[t]
  \caption{LiDAR-Camera Extinsic Calibration}
  \label{alg:ext_refinement}
  \LinesNumbered
  \KwIn{Number of iteration: $I$ \ \ \ \ \ \ \ \ \ \ \ \ \ \ \ \ \ \ \ \ \ \ \ \
    Number of LiDAR-camera data pairs: $N$
    Features: $\{\bm{n}^{c}\}$, $\{\mathcal{P}^{c}_{edge}\}$,
    $\{\mathcal{P}^{l}_{pl}\}$, $\{\mathcal{P}^{l}_{edge}\}$\\}
  \KwOut{Estimated extrinsics $\bm{T}^{c}_{l}$;}
  Initialize extrinsics $\bm{T}^{c}_{l,ini}=
    \textit{QPEP-PtPL}(\{\bm{n}^{c}\},\{\mathcal{P}^{l}_{pl}\})$\\
  \ForEach{$M\in\{1,2, \dots, N\}$}
  {
  The set of all candidate extrinsics: $\mathcal{T}=\emptyset$\\
  \While{count of iteration $<I$}
  {
  Randomly select $M$ data pairs: $\{\bm{n}^{c}\}^{\#}$, $\{\mathcal{P}^{c}_{edge}\}^{\#}$,
  $\{\mathcal{P}^{l}_{pl}\}^{\#}$, $\{\mathcal{P}^{l}_{edge}\}^{\#}$\\
  Find corresponding edge points of $\{\mathcal{P}^{l}_{edge}\}^{\#}$:
  $\{\mathcal{E}^{c}\}\subseteq\{\mathcal{P}^{c}_{edge}\}^{\#}$ given $\bm{T}^{c}_{l,ini}$\\
  Optimize extrinsics: $\bm{T}=$
  $\textit{QPEP-PtPL}(\{\bm{n}^{c}\}^{\#},$
    $\{\mathcal{P}^{l}_{pl}\}^{\#},\{\mathcal{E}^{c}\}, \{\mathcal{P}^{l}_{edge}\}^{\#})$\\
  $\mathcal{T}=\mathcal{T}\cup\bm{T}$\\
  }
  Compute the mean of $\mathcal{T}$: $\bar{\bm{T}}$\\
  $\bm{T}^{c}_{l}=\bar{\bm{T}}$ if the mean geometric error is smaller\\
  }
\end{algorithm}

\begin{figure}[t]
  \centering
  \subfigure[]
  {\label{fig:exp_rlfs01_cloud_align}\centering\includegraphics[width=0.257\textwidth]{figure/experiment/version2/rlfs01_aligned-crop.pdf}}
  \subfigure[]
  {\label{fig:exp_rles03_cloud_align}\centering\includegraphics[width=0.22\textwidth]{figure/experiment/version2/rles03_aligned-crop.pdf}}
  \caption{The LiDAR's planar points and edge points are aligned with the board plane of images with the estimated extrinsiscs.}
  \label{fig:exp_cloud_align}
\end{figure}


\subsubsection{Image Reconstruction}
One of the fundamental building blocks for camera calibration is the detection of checkerboard corners.
However, these corner detectors originally designed for frame images are not directly applicable to events due to their intrinsically asynchronous and sparse nature.
Inspired by the work done by Muglikar \textit{et al.} \cite{muglikar2021calibrate}, we resort to a learning-based method called E2VID \cite{rebecq2019high} to reconstruct high-quality frame images from the asynchronous and sparse event stream.
This method encodes events in a spatio-temporal voxel grid and uses a recurrent convolutional neural network based on the UNet architecture \cite{ronneberger2015u} to process events.
The network is already pretrained using a large number of simulated event sequences.
After the reconstruction, we can directly apply the approach in Section \ref{sec:methodology_fsc} on these images for calibration.

The reconstruction procedure is summarized in three steps:
1) Move the checkerboard before the event camera to trigger events.
2) Divide events into chunks of constant time duration ($50ms$ in our experiments). But the time duration of these chunks does not have to be constant. One could choose to define the chunks by the number of events.
3) Reconstruct images from events using E2VID. The corresponding LiDAR frame specifies the reconstruction timestamp.

Besides the above approach, we also have two possible solutions to detect corners from events, but they present several limitations.
The first solution is to aggregate events within a local spatial-temporal window to create an event map that is also an image-type representation \cite{jiao2021comparing}.
But event maps do not record intensity and contain several noisy pixels, making the traditional corner detectors inaccurate.
Another approach is to detect corners from pure events.
However, existing event-based corner detectors \cite{manderscheid2019speed} are not specifically designed for checkerboard corners, and most of them are not publicly released, inducing difficulties in calibration.
In contrast, the proposed reconstruction-based calibration scheme is easy to implement and does not require much parameter tuning.
Experimental results have demonstrated the effectiveness of the proposed reconstruction-based calibration method.


\section{Extrinsic Calibration with Multi-Frame Measurements}
\label{sec:ext_calibration}

\begin{figure}[t]
  \centering
  \subfigure[LiDAR-Frame camera sensor]
  {\label{fig:exp_real_device_1}\centering\includegraphics[width=0.21\textwidth]{figure/experiment/real_device_1_compressed.png}}
  \subfigure[LiDAR-Event camera sensor]
  {\label{fig:exp_real_device_2}\centering\includegraphics[width=0.219\textwidth]{figure/experiment/real_device_2_compressed.png}}
  \caption{The RLFS and RLES devices for calibration tests.}
  \label{fig:exp_real_device}
\end{figure}

\begin{figure}[t]
  \centering
  \subfigure[]
  {\label{fig:exp_simu_camera}\centering\includegraphics[width=0.210\textwidth]{figure/experiment/simu/simu_camera}}
  \
  \subfigure[]
  {\label{fig:exp_simu_lidar}\centering\includegraphics[width=0.21\textwidth]{figure/experiment/simu/simu_lidar}}
  \caption{Simulated sensor data: (a) an image and (b) a point cloud.}
  \label{fig:exp_simu_sensor}
\end{figure}


%%%%%%%%%%%%%%%%%%%%%%%%%%%%%%%%%%%%%%%%%%%%%%%%%%%%%%
Section \ref{sec:measure_process} has explained how checkerboard's normal vectors: $\bm{n}^{c}$, $\bm{n}^{e}$, planar points: $\mathcal{P}^{l}_{pl}$, and edge points: $\mathcal{P}^{l}_{edge}$, $\mathcal{P}^{c}_{edge}$, $\mathcal{P}^{e}_{edge}$ are extracted from each frame of LiDARs, frame cameras, and event cameras.
This section introduces how to use these features to constrain extrinsics.
We take the LiDAR-frame camera extrinsic calibration as an example without loss of generality.
At first, we can easily construct a set of point-plane pairs.
But the point-edge association is ambiguous without the initial extrinsics.
In other words, we cannot determine which boundaries a point belongs to since the board's shape is symmetric.
Thus, we follow the \textit{initialization-refinement} philosophy \cite{jiao2021robust} to estimate extrinsics in a coarse-fine manner, as summarized in Algorithm \ref{alg:ext_refinement}.

The first phase initializes the extrinsics $\bm{T}^{c}_{l,ini}$
(i.e., transformation from $()^{c}$ to $()^{l}$).
We formulate the point-to-plane optimization objective \eqref{equ:qpep-ptpl}
by utilizing planar features: $\bm{n}^{c}$ and $\mathcal{P}^{l}_{pl}$
from several frames of measurements.
The \textit{QPEP-PtPL} algorithm is applied to solve the problem.

The second phase further refines the initial solution by jointly utilizing the planar and edge points that offer stronger geometric constraints.
This phase is executed with multiple computations.
Specifically, three steps are done at line $6$-$7$:
\begin{enumerate}
  \item With $\bm{T}^{c}_{l,ini}$, we first transform $\mathcal{P}^{l}_{edge}$ into the camera frame. For each point from transformed point set, we find the nearest point in $\mathcal{P}^{c}_{edge}$ as its corresponding edge point, and then construct the point-to-line objective.
  \item From \eqref{equ:qpep-ptpl}, we propagate the covariance of each point-to-plane distance as
        $\sigma^{2}=(\bm{R}^{c}_{l}\bm{p}_{i}^{l}+\bm{t}^{c}_{l})^{\top}\bm{\Sigma}_{n^{c}}(\bm{R}^{c}_{l}\bm{p}_{i}^{l}+\bm{t}^{c}_{l})$. We filter out $10\%$ error terms if their variances are large.
  \item The overall optimization objective for the refinement is defined as the sum of all weighted point-to-plane and point-to-line residuals as
        \begin{equation}
          \label{equ:obj-refinement}
          \mathcal{L}(\bm{R}^{c}_{l},\bm{t}^{c}_{l})
          =
          w_{ptpl}\mathcal{L}_{ptpl}(\bm{R}^{c}_{l},\bm{t}^{c}_{l})
          +
          w_{ptl}\mathcal{L}_{ptl}(\bm{R}^{c}_{l},\bm{t}^{c}_{l}),
        \end{equation}
        where $w_{ptpl}$ and $w_{ptl}$ are tunned parameters which balance these two constraints.
        Similar to initialization, we can also use the \textit{QPEP-PtPL} solver.
\end{enumerate}

After getting a set of candidate extrinsics, we select one with the minimum geometric error as our resulting extrinsics.
Two examples are shown in Fig. \ref{fig:exp_cloud_align}, where LiDAR's planar points and edge points are aligned with the board plane of images with the estimated extrinsiscs.

\section{Experimental Results}
\label{sec:experiments}
\subsection{Training Details}
\cite{Kalantari2017DeepHD} provides the first dataset specifically designed for multi-exposure HDR fusion under large motion. It consists of 74 training sets, which we use to supervise the training of our model. We crop the input images to patches of size \(256 \times 256\) at a step size of 64. This totally generates 20128 training samples. To augment training samples, we randomly rotate and flip the training images. The training adopts Adam optimizer. The learning rate is initialized to \(10^{-4}\) and is reduced to \(10^{-5}\) after 20 epochs. It is observed that 40 epochs are sufficient for the training to converge.    

\subsection{Numerical Evaluation}
We numerically measure the performance of our method on the 15 test sets of \cite{Kalantari2017DeepHD}, by Peak Signal-to-Noise Ratio (PSNR) and Structure Similarity, computed in both tonemapping domain (-\(\mu\)) and HDR linear domain (-L). Visual difference metric HDR-VDP-2 is also adopted, where the parameters are set as same as in previous works \cite{wu2018end} and \cite{niu2021hdrgan}. 

Table \ref{table_metrics} compares our model with state-of-the-art models. For \cite{yan2020nonlocal} and \cite{xiong2021hierarchical}, we use the results reported in the publications. Note that \cite{sen2012robust} and \cite{hu2013hdr} are not machine learning based methods. Moreover,  \cite{Kalantari2017DeepHD} and \cite{wu2018end} apply optical flow and homography transformation to preprocess the input images respectively, and hence entail extra computation. 

Table \ref{table_metrics} shows that our method outperforms competing method in terms of PSNR-L, SSIM-$\mu$, SSIM-L and HDR-VDP-2. It ranks the second best in PSNR-$\mu$, being slightly (0.1dB) inferior to \cite{xiong2021hierarchical}. Note that \cite{xiong2021hierarchical} utilizes a pretrained model to detect ghosting regions for training, whereas our method does not require any pretrained model. The high PSNR and SSIM scores varify that our model has strong HDR reconstruction ability and can accurately restore the radiance and structure of the scene in both tonemapping domain and HDR linear domain. Furthermore, its high performance in term of HDR-VDP-2\cite{mantiuk2011hdr} performance indicates that our method can generate HDR image visually close to the target image.

\begin{table*}[ht]
\centering
\begin{tabular}{l|c|c|c|c|c}
\hline
& PSNR-$\mu$ & PSNR-L & SSIM-$\mu$ & SSIM-L & HDR-VDP-2 \\
\hline
\bfseries Sen & 40.97 & 38.36 & 0.9830 & 0.9746 & 60.60\\
\hline
\bfseries Hu  & 35.65 & 30.80 & 0.9725 & 0.9491 & 58.34\\
\hline
\bfseries Kalantari & 42.69 & 41.22 & 0.9888 & 0.9845 & 65.05\\
\hline
\bfseries DeepHDR& 41.99 & 41.22 & 0.9878 & 0.9859 & \underline{65.91}\\
\hline
\bfseries AHDR & 43.62 & 41.03 & 0.9900  &\underline{0.9883} & 63.85 \\
\hline 
\bfseries NHDRRNet& 42.414 & - & 0.9887 & - & 61.21 \\
\hline 
\bfseries HDR-GAN &43.92 & \underline{41.57} &\underline{0.9905} &0.9865 & 65.45\\
\hline 
\bfseries HFNet & \textbf{44.28} & 41.47 & - & - & - \\
\hline 
\bfseries Ours & \underline{44.18} & \textbf{42.19}&\textbf{0.9912} & \textbf{0.9883}& \textbf{67.07} \\
\hline
\end{tabular}
\caption{Numerical performance of the proposed model, evaluated on the dataset by Kalantari-Ramamoorthi. The best and second best results for each metric are marked in \textbf{bold} and \underline{underlined}, respectively}
\label{table_metrics}
\end{table*}

\subsection{Visual Performance Evaluation}

\begin{figure*}[!htb]
\centering
\includegraphics[width=\textwidth]{experiments/kalantari_test.png}
\caption{Visual comparison on the test set of Kalantari-Ramamoorthi dataset. Zoom-in views of reconstruction by each method are presented on the saturated regions that contain moving objects. Our network built with gated Swin Transformer yields noticeably better visual results than other methods.}
\label{fig_kalantari_test}
\end{figure*}
Fig. \ref{fig_kalantari_test} present the visual performance of our method and comparable methods on two examples from \cite{Kalantari2017DeepHD}. We present the zoom-in views of two challenging cases, where large saturated regions contain substantial non-rigid motion in the reference image. The two patch-based methods do not reconstruct the missing details in the saturated regions, as they heavily rely on the details provided by the reference image for registration. Image reconstructed by the optical flow based method \cite{Kalantari2017DeepHD} suffers motion blur artifacts. This is because the convolutions of DeepHDR and HDR-GAN have limited receptive fields, and hence are hampered to repair missing content in misaligned regions by aligned regions. The gating mechanism of AHDR is only applied to low-level features, so the high-level outliers may deteriorate the HDR fusion. In contrast to comparable methods, our model remarkably overcomes the ghosting artifacts.

\begin{figure}[ht]
\centering
\includegraphics[width=\columnwidth]{experiments/sen_test.pdf}
\caption{Visual performance comparison on example images from the dataset by Sen et al. Zoom in views on challenging areas are presented. Although the ground truth is unavailable, it can be clearly observed that our method visually performs better than comparable methods.}
\label{sen_test}
\end{figure}

\begin{figure}[ht]
\centering
\includegraphics[width=\columnwidth]{experiments/tursun_test.pdf}
\caption{Visual performance comparison on example images from the dataset by Tursun et al. Compared to state of the art methods, our method suffers less ghosting artifact.}
\label{tursun_test}
\end{figure}

Fig.\ref{sen_test} and Fig.\ref{tursun_test} present visual performance of our method on two examples from benchmark datasets \cite{sen2012robust} and \cite{tursun2016objective}. As these test datasets   do not provide ground truth image. we mark the visual difference on the results generated by different methods. It can be seen that our method suffers less artifacts than other methods in various scenes with various motion patterns, achieving better visual results. Our method creates high-quality HDR more robustly and generalizes well. 

\subsection{Ablation Study}

\begin{table}[h]
\centering
\resizebox{\columnwidth}{!}{
\begin{tabular}{l|c|c|c|c|c}
\hline
                         & PSNR-$\mu$ & PSNR-l & SSIM-$\mu$ & SSIM-l & HDR-VDP-2 \\ \hline
restormer(w/o ssim loss) & 44.00  & 41.5   & 0.9906 & 0.9873 & 64.72  \\ \hline
Ours(w/o ssim loss)      & 44.07  & 41.83  & 0.9909 & 0.9879 &  64.78  \\ \hline
Ours                     & 44.18  & 42.19  & 0.9912 & 0.9883 & 67.07      \\ \hline
\end{tabular}
}
\caption{Experimental results of ablation study. We compare using Gated Swin Transformer v.s. Gated Transformer, and the combined loss function v.s. the traditional $l_{1}$ norm loss function.}
\label{table_ablation_block_loss}
\end{table}

We verify various components of our method, including Swin Transformer, loss function, and gating mechanism by ablation study.

\subsubsection{Ablation Study on Block Design}
Our model has similar architecture to Restormer, which uses modified Transformer, whereas we use modified Swin Transformer as the building unit. For comparison, we replace the residual modules in each block in our model with multiple transformer layers as in Restormer, with same number of transformer layers. Table \ref{table_ablation_block_loss} presents the results, which show that using Swin Transformer achieves superior performance in all measures. The reason is that the attention module of Restormer is computed channel-wise, but forgoes the cross-exposure spatial dependency to repair the non-aligned area. 

\subsubsection{Ablation Study on Loss Function}
We trained our model under different loss function configurations, as shown in \ref{table_ablation_block_loss}. The results validate that the SSIM loss benefits detail reconstruction.

\subsubsection{Ablation Study on Gating Mechanism}
\begin{table}[h]
\resizebox{\columnwidth}{!}{
\begin{tabular}{l|c|c|c|c|c}
\hline
           & PSNR-$\mu$ & PSNR-l & SSIM-$\mu$ & SSIM-l & HDR-VDP-2 \\ \hline
w/o gating & 43.14  & 41.03  & 0.9904 & 0.9868 &     64.88      \\ \hline
one gating & 43.44  & 41.42  & 0.9909 & 0.9882 &    67.13   \\ \hline
Ours       & 43.61  & 41.74  & 0.9909 & 0.9881 & 66.96     \\ \hline
\end{tabular}
}
\caption{Ablation experimental results to verify the effectiveness of the gating mechanism}
\label{table_ablation_gating}
\end{table}

The gating mechanism is an important component in our model. Ablation study is conducted in the gating mechanism as follows.

\textbf{w/o gating}: The gating mechanism is not used in the feed forward network of all transformer layers in the model, that it, our GST unit degenerate to the vanilla Swin Transformer.

\textbf{one gating}: The gating mechanism is only used in the first Swin Transformer layers subsequent to the embedding layer, but not used for other layers. 

 Table \ref{table_ablation_gating} shows the results of the ablation experiments, where the model is trained for 20 epochs. By removing the gating mechanism, the network relies on self-attention for image alignment, resulting in the lowest performance. On top of it, adding gates to low level layers notably improves the HDR reconstruction. Furthermore, by integrating the gating mechanism with all Swin Transformer layers, the model effectively inpaints information in non-aligned regions and obtains the highest HDR reconstruction results, thus validates the effectiveness of the gating mechanism in our model.

\begin{figure}[t]
  \centering\includegraphics[width=0.45\textwidth]{figure/experiment/version3/err_simu_data_2_methods-crop}
  \caption{Calibration error of the proposed method and \textit{Zhou-MATLAB} on calibration data at different levels of noise. The proposed method can achieve calibration accracy (the top three rows) ranging from $[0.068,0.400]deg$, $[0.0004, 0.012]m$, $[0.0003, 0.024]m$, $[0.002, 0.023]m$ at    $EGT_{\bm{R}}$, $EGT_{\bm{t}}$, and $MGE$ respectively. Our method achieves the plane fitting accuracy (the fourth row) that is computed as the relative angle with the ground-truth normal ranging from $[0.0104,12.8701]deg$.}
  \label{fig:exp_err_plot_simu_data}
\end{figure}

\begin{figure*}[t]
  \centering
  \subfigure[Median rotation error]
  {\label{fig:exp_rlfs_error_rot}\centering\includegraphics[width=0.232\textheight]{figure/experiment/version2/rlfs-rot-error-crop.pdf}}
  \subfigure[Median translation error]
  {\label{fig:exp_rlfs_error_tsl}\centering\includegraphics[width=0.232\textheight]{figure/experiment/version2/rlfs-tsl-error-crop.pdf}}
  \subfigure[Median MGE]
  {\label{fig:exp_rlfs_error_mge}\centering\includegraphics[width=0.263\textheight]{figure/experiment/version2/rlfs-mg-error-crop.pdf}}
  \caption{Median rotation error, translation error, and MGE of ours and \textit{Zhou-MATLAB} on calibration data: \textit{RLFS01}-\textit{RLFS03}. The median errors are computed using the set of candidate extrinsics at line $8$ of Algorithm \ref{alg:ext_refinement}.}
  \label{fig:exp_rlfs_error}
\end{figure*}

\begin{table*}[]
  \centering
  \caption{Calibration results on \textit{RLFS}.
    $\downarrow$/$\uparrow$ indicates that the lower/higher the value, the better the score.}
  \renewcommand\arraystretch{0.9}
  \renewcommand\tabcolsep{4.6pt}
  \scriptsize
  \begin{tabular}{ccccccccccccc}
    \toprule
    \multirow{2}{*}{Dataset}                              &
    \multirow{2}{*}{Method}                               &
    \multicolumn{4}{c}{Quaternion}                        &
    \multicolumn{3}{c}{Translation {$[m]$}}               &
    \multirow{2}{*}{$EGT_{\mathbf{R}}\ [deg,\downarrow]$} &
    \multirow{2}{*}{$EGT_{\mathbf{t}}\ [m,\downarrow]$}   &
    \multirow{2}{*}{\textit{MGE}$\ [m,\downarrow]$}       &
    \multirow{2}{*}{$N_{\text{detect}}/N\ [\uparrow]$}                                                                                                                                                                                                                                                     \\
    \cline{3-9}
                                                          &                         & \multicolumn{1}{c}{$w$} & \multicolumn{1}{c}{$x$} & \multicolumn{1}{c}{$y$} & \multicolumn{1}{c}{$z$}
                                                          & \multicolumn{1}{c}{$x$} & \multicolumn{1}{c}{$y$} & \multicolumn{1}{c}{$z$} &                         &                         &                &                                                                                             \\
    \midrule[0.03cm]

                                                          & \textit{CAD}            & $\dgt{0.664}$           & $\dgt{0.664}$           & $\dgt{-0.242}$          & $\dgt{0.242}$           & $\dgt{0.000}$  & $\dgt{-0.059}$ & $\dgt{-0.070}$
                                                          & $\dgt{-}$               & $\dgt{-}$               & $\dgt{-}$               & $\dgt{-}$                                                                                                                                                        \\
                                                          & \textit{Fake GT}        & $\dgt{0.666}$           & $\dgt{0.667}$           & $\dgt{-0.233}$          & $\dgt{0.238}$           & $\dgt{-0.008}$ & $\dgt{-0.069}$ & $\dgt{-0.084}$
                                                          & $\dgt{-}$               & $\dgt{-}$               & $\dgt{-}$               & $\dgt{-}$                                                                                                                                                        \\
    \midrule

    \multirow{2}{*}{\textit{RLFS01}}
                                                          & \textit{Ours}           & $0.667$                 & $0.666$                 & $-0.233$                & $0.239$                 & $-0.003$       & $-0.074$       & $-0.091$       & $\bm{0.130}$ & $\bm{0.009}$ & $\bm{0.009}$ & $\bm{17}/18$ \\
                                                          & \textit{Zhou-MATLAB}    & $0.669$                 & $0.662$                 & $-0.238$                & $0.240$                 & $0.016$        & $-0.071$       & $-0.093$       & $0.850$      & $0.026$      & $0.011$      & $10/18$      \\
    \midrule

    \multirow{2}{*}{\textit{RLFS02}}
                                                          & \textit{Ours}           & $0.665$                 & $0.668$                 & $-0.234$                & $0.238$                 & $-0.000$       & $-0.060$       & $-0.090$       & $\bm{0.293}$ & $\bm{0.014}$ & $\bm{0.007}$ & $\bm{26}/29$ \\
                                                          & \textit{Zhou-MATLAB}    & $0.670$                 & $0.663$                 & $-0.227$                & $0.244$                 & $-0.008$       & $-0.110$       & $-0.095$       &
    $1.223$                                               & $0.042$                 & $0.013$                 & $21/29$                                                                                                                                                                                    \\
    \midrule

    \multirow{2}{*}{\textit{RLFS03}}
                                                          & \textit{Ours}           & $0.667$                 & $0.667$                 & $-0.233$                & $0.237$                 & $-0.010$       & $-0.068$       & $-0.092$       & $\bm{0.136}$ & $\bm{0.008}$ & $\bm{0.008}$ & $\bm{15}/15$ \\
                                                          & \textit{Zhou-MATLAB}    & $0.669$                 & $0.661$                 & $-0.238$                & $0.244$                 & $0.029$        & $-0.077$       & $-0.100$       & $1.120$      & $0.041$      & $0.013$      & $13/15$      \\

    \bottomrule[0.03cm]
  \end{tabular}
  \label{tab:rlfs_calibration_results}
\end{table*}

\begin{figure*}[t]
  \centering
  \subfigure[\textit{RLFS01}]
  {\label{fig:exp_rlfs_calibration_image_01}\centering\includegraphics[width=0.243\textwidth]{figure/experiment/version2/rlfs01-crop_compressed}}
  \subfigure[\textit{RLFS02}]
  {\label{fig:exp_rlfs_calibration_image_02}\centering\includegraphics[width=0.241\textwidth]{figure/experiment/version2/rlfs02-crop_compressed}}
  \subfigure[\textit{RLFS03}]
  {\label{fig:exp_rlfs_calibration_image_03}\centering\includegraphics[width=0.220\textwidth]{figure/experiment/version2/rlfs03-crop_compressed}}
  \subfigure[\textit{RLFS02} (cloud)]
  {\label{fig:exp_rlfs_calibration_cloud_02_cloud}\centering\includegraphics[width=0.260\textwidth]{figure/experiment/version2/cloud_rlfs02-crop}}
  \caption{Back-projection of the LiDAR points om images using the extrinsics from Table \ref{tab:rlfs_calibration_results} from our algorithm (red) and \textit{Zhou-MATLAB} (green).}
  \label{fig:exp_rlfs_calibration_image}
\end{figure*}


\subsection{Implementation Details}
\label{sec:exp_implementation}

We implement the proposed LCE-Calib in MATLAB.
We empirically set $\mu_1=0.9$ and $\mu_2=1$ in the checkboard point selection,
and $I=100$ in Algorithm \ref{alg:ext_refinement}.
Hyperparameters of the objective \eqref{equ:obj-refinement}: $w_{ptpl}$ and $w_{ptl}$, should be set according to a specific sensor configuration,
affected by the LiDAR's density and camera's resolution.

\subsubsection{Sensor Suites}
The proposed method is tested on both the simulated and real-world sensor suite.
The LiDAR and camera on real platforms are hardware-synchronized.
Hyperparameters: the number of checkerboard's inner patterns and board size should be set for a specific calibration target.
\begin{enumerate}
  \item \textit{The Simulated Sensor Suite} (SSS) is built on the Gazebo \cite{koenig2004design}. It consists of a simulated 16-beam LiDAR and frame camera with the same specification to RLFS. The ground-truth intrinsics and extrinsics are provided.

  \item \textit{The Real-World LiDAR-Frame Camera Sensor} (RLFS) contains a VLP-16 LiDAR\footnote{\url{https://velodynelidar.com/products/puck}} and a high-resolution camera (GSML-AR0143-H090, resolution: $1280$H$\times 720$V)\footnote{\url{https://docs.miivii.com/product/apex/xavier/manual/en/common/05.EN_gmsl.html}}, as shown in Fig. \ref{fig:exp_real_device_1}.

  \item \textit{The Real-World LiDAR-Event Camera Sensor} (RLES) is made for the mapping application and shown in Fig. \ref{fig:exp_real_device_2}. It is installed with a VLP-16 LiDAR and an event camera (DAVIS346, resolution: $360$H$\times 240$V)\footnote{\url{https://inivation.com/wp-content/uploads/2019/08/DAVIS346.pdf}}.
\end{enumerate}

Both frame and event cameras use the pinhole projection model.
And the DAVIS346 event camera outputs RGB frames that are used to verify the frame camera-based calibration.

\subsubsection{Evaluation Metrics}
We introduce two metrics to assess the LiDAR-camera calibration results from different aspects.
\begin{enumerate}
  \item \textit{Error Compared With Ground Truth (EGT)} computes the distance between the ground truth and the estimated values in terms of rotation and translation as
        \begin{equation}
          \begin{aligned}
            EGT_{\bm{R}}
             & =
            ||\ln(\bm{R}_{gt}\bm{R}_{est}^{-1})^{\vee}||, \\
            EGT_{\bm{t}}
             & =
            ||\bm{t}_{gt}-\bm{t}_{est}||.
          \end{aligned}
        \end{equation}
  \item \textit{Mean Geometric Error (MGE)} is the sum of the mean planar error (MPE) and mean edge error (MEE). MPE and MEE compute the mean point-to-plane and point-to-line distance between the 3D checkerboard plane in the camera frame and checkerboard points in the LiDAR frame respectively.
\end{enumerate}

The perfect ground truth is unknown in real-world experiments.
We first estimate 3D corner points from both images and point clouds.
We then manually check and modify them.
With 3D-3D point correspondences that are found, we can compute the ``fake'' ground truth with the ICP algorithm.
Our proposed method is compared with a baseline method (denoted by \textit{Zhou-MATLAB}) that has been released in MATLAB\footnote{\url{https://www.mathworks.com/help/lidar/ug/lidar-and-camera-calibration.html}},
This method implements the SOTA LiDAR-camera calibration method \cite{zhou2018automatic} but formulates a different objective function.

\begin{figure*}[t]
  \centering
  \subfigure[Median rotation error]
  {\label{fig:exp_rles_error_rot}\centering\includegraphics[width=0.31\textwidth]{figure/experiment/version2/rles-rot-err-crop.pdf}}
  \subfigure[Median translation error]
  {\label{fig:exp_rles_error_tsl}\centering\includegraphics[width=0.31\textwidth]{figure/experiment/version2/rles-tsl-err-crop.pdf}}
  \subfigure[Median MGE]
  {\label{fig:exp_rles_error_tsl}\centering\includegraphics[width=0.342\textwidth]{figure/experiment/version2/rles-mg-err-crop.pdf}}
  \caption{Median rotation error, translation error, and MGE at multiple runs	of our method on calibration data: \textit{RLES01-LF}, \textit{RLES01-LE}, \textit{RLES02-LF}, \textit{RLES02-LE}, \textit{RLES03-LF}, and \textit{RLES03-LE}. The median errors are computed using the set of candidate extrinsics at line $8$ of Algorithm \ref{alg:ext_refinement}.}
  \label{fig:exp_rles_error}
\end{figure*}

\begin{table*}[]
  \centering
  \caption{Calibration results on \textit{RLES}. $\downarrow$/$\uparrow$ indicates that the lower/higher the value, the better the score.}
  \renewcommand\arraystretch{0.9}
  \renewcommand\tabcolsep{4.6pt}
  \scriptsize
  \begin{tabular}{ccccccccccccc}
    \toprule
    \multirow{2}{*}{Dataset}                               &
    \multirow{2}{*}{Method}                                &
    \multicolumn{4}{c}{Quaternion}                         &
    \multicolumn{3}{c}{Translation {$[m]$}}                &
    \multirow{2}{*}{$EGT_{\mathbf{R}}\ [deg, \downarrow]$} &
    \multirow{2}{*}{$EGT_{\mathbf{t}}\ [m, \downarrow]$}   &
    \multirow{2}{*}{\textit{MGE}$\ [m, \downarrow]$}       &
    \multirow{2}{*}{$N_{\text{detect}}/N\ [\uparrow]$}                                                                                                                                                                                                                                       \\
    \cline{3-9}
                                                           &                         & \multicolumn{1}{c}{$w$} & \multicolumn{1}{c}{$x$} & \multicolumn{1}{c}{$y$} & \multicolumn{1}{c}{$z$}
                                                           & \multicolumn{1}{c}{$x$} & \multicolumn{1}{c}{$y$} & \multicolumn{1}{c}{$z$} &                         &                         &                                                                                               \\
    \midrule[0.03cm]

                                                           & \textit{CAD}            & $\dgt{0.500}$           & $\dgt{0.500}$           & $\dgt{-0.500}$          & $\dgt{0.500}$           & $\dgt{0.126}$  & $\dgt{-0.065}$ & $\dgt{-0.080}$
                                                           & $-$                     & $-$                     & $-$                     & $-$                                                                                                                                               \\
    \midrule
    \midrule

                                                           & \textit{Fake GT}        &
    $\dgt{0.502}$                                          & $\dgt{0.504}$           & $\dgt{-0.498}$          & $\dgt{0.495}$           & $\dgt{0.112}$           & $\dgt{-0.094}$          & $\dgt{-0.108}$
                                                           & $-$                     & $-$                     & $-$                     & $-$                                                                                                                                               \\
    \midrule

    \multirow{2}{*}{\textit{RLES01-LF}}
                                                           & \textit{Ours}
                                                           & $0.499$                 & $0.505$                 & $-0.498$                & $0.498$                 & $0.126$                 & $-0.095$       & $-0.106$       & $\bm{0.629}$   & $\bm{0.032}$ & $0.037$      & $\bm{20}/20$
    \\
                                                           & \textit{Zhou-MATLAB}
                                                           & $0.502$                 & $0.498$                 & $-0.502$                & $0.498$                 & $0.132$                 & $-0.094$       & $-0.104$       & $1.151$        & $0.038$      & $\bm{0.036}$ & $\bm{20}/20$
    \\
    \midrule

    \multirow{2}{*}{\textit{RLES01-LE}}
                                                           & \textit{Ours}
                                                           & $0.497$                 & $0.501$                 & $-0.501$                & $0.500$                 & $0.150$                 & $-0.092$       & $-0.096$       & $\bm{0.839}$   & $\bm{0.041}$ & $0.024$      & $\bm{33}/36$
    \\
                                                           & \textit{Zhou-MATLAB}
                                                           & $0.496$                 & $0.500$                 & $-0.500$                & $0.503$                 & $0.161$                 & $-0.097$       & $-0.100$       & $1.258$        & $0.050$      & $\bm{0.023}$ & $\bm{33}/36$
    \\
    \midrule
    \midrule

                                                           & \textit{Fake GT}        &
    $\dgt{0.500}$                                          & $\dgt{0.505}$           & $\dgt{-0.498}$          & $\dgt{0.497}$           & $\dgt{0.129}$           & $\dgt{-0.080}$          & $\dgt{-0.109}$
                                                           & $-$                     & $-$                     & $-$                                                                                                                                                                         \\
    \midrule

    \multirow{2}{*}{\textit{RLES02-LF}}
                                                           & \textit{Ours}
                                                           & $0.497$                 & $0.506$                 & $-0.502$                & $0.495$                 & $0.131$                 & $-0.064$       & $-0.102$       & $\bm{0.677}$   & $\bm{0.018}$ & $\bm{0.031}$ & $\bm{34}/36$
    \\
                                                           & \textit{Zhou-MATLAB}
                                                           & $0.495$                 & $0.504$                 & $-0.504$                & $0.497$                 & $0.160$                 & $-0.060$       & $-0.090$       & $0.927$        & $0.041$      & $0.037$      & $32/36$
    \\
    \midrule

    \multirow{2}{*}{\textit{RLES02-LE}}
                                                           & \textit{Ours}
                                                           & $0.498$                 & $0.507$                 & $-0.499$                & $0.496$                 & $0.130$                 & $-0.071$       & $-0.107$       & $\bm{0.331}$   & $\bm{0.009}$ & $\bm{0.028}$ & $\bm{34}/36$
    \\
                                                           & \textit{Zhou-MATLAB}
                                                           & $0.491$                 & $0.505$                 & $-0.505$                & $0.498$                 & $0.152$                 & $-0.049$       & $-0.096$       & $1.270$        & $0.041$      & $0.034$      & $31/36$
    \\
    \midrule

    \multirow{2}{*}{\textit{RLES03-LF}}
                                                           & \textit{Ours}
                                                           & $0.497$                 & $0.511$                 & $-0.496$                & $0.496$                 & $0.116$                 & $-0.055$       & $-0.110$       & $\bm{0.781}$   & $\bm{0.028}$ & $0.027$      & $\bm{29}/36$
    \\
                                                           & \textit{Zhou-MATLAB}
                                                           & $0.508$                 & $0.504$                 & $-0.494$                & $0.495$                 & $0.100$                 & $-0.104$       & $-0.114$       & $1.034$        & $0.039$      & $\bm{0.024}$ & $26/36$
    \\
    \midrule

    \multirow{2}{*}{\textit{RLES03-LE}}
                                                           & \textit{Ours}
                                                           & $0.504$                 & $0.503$                 & $-0.498$                & $0.495$                 & $0.119$                 & $-0.090$       & $-0.108$       & $0.554$        & $\bm{0.014}$ & $0.030$      & $\bm{29}/36$
    \\
                                                           & \textit{Zhou-MATLAB}
                                                           & $0.497$                 & $0.507$                 & $-0.499$                & $0.497$                 & $0.124$                 & $-0.066$       & $-0.101$       & $\bm{0.401}$   & $0.016$      & $\bm{0.028}$ & $27/36$
    \\
    \bottomrule[0.03cm]
  \end{tabular}
  \label{tab:rles_calibration_results}
\end{table*}


\subsection{Calibration Results on Simulated Data}
\label{sec:exp_calib_simulate}
We first verify the LiDAR-frame camera calibration with synthetic data.
% We use the open-source Gazebo software \cite{koenig2004design} to build the simulated frame camera and 16-beam LiDAR.
The sampled image and point cloud are shown in Fig. \ref{fig:exp_simu_sensor}.
The LiDAR and camera are given sufficient rotational offset: $[-70,-70,150]deg$ at roll, pitch and yaw respectively, as well as translational offset: $[0.4185, -0.3050, -0.1476]m$ at $x-$, $y-$, and $z-$ axes respectively.

We place the checkerboard in front of the sensor with sufficient rotation and translation and obtain
$\bm{30}$ pairs of noise-free LiDAR-camera data.
To test the robustness and limitations of our method, we further generate another $\bm{9}$ sets of data.
First, we add zero-mean Gaussian noise on depth measurements with an increasing standard deviation of
$\sigma\in\{0.8, 1.6, 2.4, 3.2, 4.0, 4.8, 6.0, 8.0, 10.0\}cm$ to the noise-free LiDAR data.
Second, we additionally process points which are next to the checkerboard's boundaries by adding another Gaussian noise.
This is because we want to mimic real-world data since the practical emmitted laser has a divergence angle rather than a strict line \cite{yuan2021pixel}.

In Fig. \ref{fig:exp_err_plot_simu_data}, we plot the calibration error.
Benefiting from the proposed feature extraction and globally optimization, our method overall has more accurate results than \textit{Zhou-MATLAB} even under large noise.
A key step in feature extraction is estimating the normal vector of the checkerboard plane from LiDARs.
Both noisy/outlier points and improper plane fitting may negatively affect the estimates.
This figure also shows the mean angular error between and ground-truth normals and estimated normals.
\textit{Zhou-MATLAB} cannot estimate accurate planar coefficients.
Regarding the optimization part, our method has rich planar and edge constraints to enforce the accuracy.
Also, the globally optimal solution in our method is essential to these results.
Moreover, we find that the unsupervised metric: \textit{MGE} is also suitable to evaluate calibration results, since its curve is consistent with the curves of \textit{EGT}.


\begin{figure}[]
  \centering
  \includegraphics[width=0.49\textwidth]{figure/experiment/rlfs/error_num_frame_rlfs02-crop.pdf}
  \caption{Calibration error along with frame number on \textit{RLFS01}.}
  \label{fig:calibration_error_frame_rlfs}
\end{figure}

\begin{figure}[t]
  \centering
  \subfigure[\textit{RLES01-LF}]
  {\label{fig:exp_rles_calibration_image_01_lf}\centering\includegraphics[width=0.22\textwidth]{figure/experiment/version2/rles01-crop}}
  \subfigure[\textit{RLES01-LE}]
  {\label{fig:exp_rles_calibration_image_01_le}\centering\includegraphics[width=0.22\textwidth]{figure/experiment/version2/rles02-crop}}
  \subfigure[\textit{RLES02-LF}]
  {\label{fig:exp_rles_calibration_image_02_lf}\centering\includegraphics[width=0.22\textwidth]{figure/experiment/version2/rles03-crop}}
  \subfigure[\textit{RLES02-LE}]
  {\label{fig:exp_rles_calibration_image_02_le}\centering\includegraphics[width=0.22\textwidth]{figure/experiment/version2/rles04-crop}}
  \subfigure[\textit{RLES03-LF}]
  {\label{fig:exp_rles_calibration_image_03_lf}\centering\includegraphics[width=0.22\textwidth]{figure/experiment/version2/rles05-crop}}
  \subfigure[\textit{RLES03-LE}]
  {\label{fig:exp_rles_calibration_image_03_le}\centering\includegraphics[width=0.22\textwidth]{figure/experiment/version2/rles06-crop}}
  \subfigure[\textit{RLES02} (cloud)]
  {\label{fig:exp_rles_calibration_02_cloud}\centering\includegraphics[width=0.22\textwidth]{figure/experiment/rles/calibrate_cloud_rles02-crop}}
  \subfigure[\textit{RLES03} (cloud)]
  {\label{fig:exp_rles_calibration_03_cloud}\centering\includegraphics[width=0.22\textwidth]{figure/experiment/rles/calibrate_cloud_rles03-crop}}
  \caption{Back-projection of the LiDAR points on images using the extrinsics in Table \ref{tab:rles_calibration_results}: our algorithm (red) and \textit{Zhou-MATLAB} (green).}
  \label{fig:exp_rles_calibration_image}
\end{figure}



\subsection{Calibration Results on Real-World Data}
\label{sec:exp_calib_realworld}

This section tests the proposed method with two real-world sensor suites.
Similarly, we move a checkerboard before sensors slowly to collect several groups of calibration data. Experimental results are shown in the following sections.

%%%%%%%%%%%%%%%%%%%%%%%%%%%%%%%%%%%%%%%%%%
\subsubsection{Extrinsic Calibration of RLFS}
\label{sec:exp_rlfs}
We first verify the calibration method with the LiDAR-frame camera setup.
We collect three groups of data called \textit{RLFS01}, \textit{RLFS02}, and \textit{RLFS03} to test our approach.
\textit{RLFS01} and \textit{RLFS02} were collected in an indoor corridor. Besides the checkerboard, other objects such as grounds, walls, and square pillars appear as planar surfaces.
\textit{RLF03} were collected in an outdoor garden. Several noisy objects such as buildings, poles, and lawns exist, affecting the checkerboard detection from point clouds.

We firstly employ the MATLAB camera calibration toolbox\footnote{\url{https://www.mathworks.com/help/vision/camera-calibration.html}} to calibrate the camera's intrinsics with these collected data and then calibrate extrinsics.
Our method successfully recovers the extrinsics.
Fig. \ref{fig:exp_rlfs_error} plots errors with increasing frames of measurements and Table \ref{tab:rlfs_calibration_results} gives quantitative sample results.
$N_{\text{detect}}/N$ shows the portion of successfully detecting the checkerboard from images or point clouds.
\textit{Zhou-MATLAB} only uses a portion of measurements since it fails to detect checkerboards from several LiDAR frames and correctly detect features, thus obtaining inaccurate extrinsics.
These results are consistent with those in Section \ref{sec:exp_calib_simulate}, where our method outperforms the baseline and reaches up to $<0.3deg$ and $<1.5cm$ calibration error as well as $<0.9cm$ \textit{MGE} in both indoor and outdoor environments.

Fig. \ref{fig:exp_rlfs_calibration_image} shows the back-projection results using extrinsics in Table \ref{tab:rlfs_calibration_results}.
Circular points indicate detected edge points.
Our method can provide better results since the projections of edge points nearly stay on the the checkerboard's boundary.
We additionally show the calibration error along with the number of frames used in
calibration. The error gradually decreases since more constraints help to enforce the accuracy. One example on \textit{RLFS01} is shown in Fig. \ref{fig:calibration_error_frame_rlfs}.

\subsubsection{Extrinsic Calibration of RLES}
\label{sec:exp_rles}

We further verify the complete LCE-Calib method with RLES.
Three groups of calibration data are collected in an indoor room under different lighting conditions, which are called \textit{RLES01}, \textit{RLES02}, and \textit{RLES03} respectively.
In contrast to the last section, we additionally compare calibration results of using frame images (\textit{X-LF}) and reconstructed images (\textit{X-LE}) from event streams, respectively.
Fig. \ref{fig:exp_rles_error} plots calibration errors of the proposed method and Table \ref{tab:rles_calibration_results} reports quantitative results.
% \textit{RLFS01} and \textit{RLFS02}/\textit{RLFS03} are collected on different days.
% We provide two ``fake'' ground-truth extrinsics for them.

The calibration errors consistently decrease with more frames.
Though, compared with Table \ref{tab:rlfs_calibration_results}, the calibration errors are larger.
We consider that this is mainly caused by the low resolution of the event camera, affecting the inner pattern detection accuracy.
Cameras' low resolution also makes \textit{MGE} values not consistent with \textit{EGT} in some cases since boards' positions are not detected very precisely.
But we can still use this metric for selecting good extrinsics.
% In most cases, the proposed method outperforms the baseline algorithm.
% The proposed method has comparable
From Fig. \ref{fig:exp_rles_error}, we observe that extrinsics converge if around $\bm{15}$ frames are given. More data can further enforce accuracy.
For event cameras, inner patterns and boundaries of the board are reconstructed clearly. Thus, the calibration accuracy does not have much degradation.
A sufficient number of frames are used in calibration, whether using frame images or images reconstructed from event streams, and the extrinsics are successfully recovered.
Fig. \ref{fig:exp_rles_calibration_image} shows the back-projection results of LiDAR points using extrinsics in Table \ref{tab:rles_calibration_results}.
The checkerboard detection results from LiDAR points are also shown.
In \textit{RLES03-LF}, extrinsics from our method can help LiDAR points better align on the board's boundaries than \textit{Zhou-MATLAB}.


\begin{table}[t]
	\centering
	\caption{Calibration results on \textit{RLFS} with our proposed method by disabling some modules. $\downarrow$/$\uparrow$ indicates that the lower/higher the value, the better the score. The first two results are marked as bold.}
	\renewcommand\arraystretch{1.0}
	\renewcommand\tabcolsep{6.0pt}
	\scriptsize
	\begin{tabular}{ccccc}
		\toprule
		\multirow{1}{*}{Dataset}                              &
		\multirow{1}{*}{Method}                               &
		\multirow{1}{*}{$EGT_{\mathbf{R}}\ [deg,\downarrow]$} &
		\multirow{1}{*}{$EGT_{\mathbf{t}}\ [m,\downarrow]$}   &
		\multirow{1}{*}{\textit{MGE}$\ [m,\downarrow]$}
		\\
		\midrule[0.03cm]

		\multirow{7}{*}{\textit{RLFS01}}

		                                                      & \textit{WO. PP}   & $0.517$      & $0.023$      & $0.015$      \\
		                                                      & \textit{WO. UM}   & $\bm{0.119}$ & $0.010$      & $\bm{0.008}$ \\
		                                                      & \textit{WO. PtPL} & $0.292$      & $0.011$      & $\bm{0.008}$ \\
		                                                      & \textit{WO. PtL}  & $0.556$      & $0.018$      & $\bm{0.008}$ \\
		                                                      & $I=1$             & $0.195$      & $\bm{0.007}$ & $0.009$      \\
		                                                      & $I=5$             & $0.262$      & $0.012$      & $\bm{0.008}$ \\
		                                                      & \textit{Ours}     & $\bm{0.130}$ & $\bm{0.009}$ & $0.009$      \\
		\midrule

		\multirow{7}{*}{\textit{RLFS02}}

		                                                      & \textit{WO. PP}   & $0.369$      & $0.021$      & $0.015$      \\
		                                                      & \textit{WO. UM}   & $0.338$      & $0.015$      & $0.008$      \\
		                                                      & \textit{WO. PtPL} & $0.314$      & $\bm{0.012}$ & $0.008$      \\
		                                                      & \textit{WO. PtL}  & $0.653$      & $0.021$      & $\bm{0.007}$ \\
		                                                      & $I=1$             & $\bm{0.277}$ & $\bm{0.014}$ & $0.008$      \\
		                                                      & $I=5$             & $0.342$      & $0.015$      & $0.008$      \\
		                                                      & \textit{Ours}     & $\bm{0.293}$ & $\bm{0.014}$ & $\bm{0.007}$ \\
		\midrule

		\multirow{7}{*}{\textit{RLFS03}}

		                                                      & \textit{WO. PP}   & $0.378$      & $0.018$      & $0.016$      \\
		                                                      & \textit{WO. UM}   & $\bm{0.207}$ & $\bm{0.007}$ & $\bm{0.008}$ \\
		                                                      & \textit{WO. PtPL} & $0.338$      & $0.016$      & $0.009$      \\
		                                                      & \textit{WO. PtL}  & $0.289$      & $0.019$      & $\bm{0.007}$ \\
		                                                      & $I=1$             & $0.475$      & $0.024$      & $\bm{0.008}$ \\
		                                                      & $I=5$             & $0.324$      & $0.017$      & $\bm{0.008}$ \\
		                                                      & \textit{Ours}     & $\bm{0.136}$ & $\bm{0.008}$ & $\bm{0.008}$ \\
		\bottomrule[0.03cm]
		\multicolumn{5}{l}{WO.: without. PP: point projection. UM: uncertainty modelling of normal vectors.}                   \\
		% \multicolumn{5}{l}{}                                                       \\
		\multicolumn{5}{l}{PtPL: point-to-plane constriants. PtL: point-to-line constriants.}                                  \\
		\multicolumn{5}{l}{$I$: the computation number in Algorithm \ref{alg:ext_refinement}.}                                 \\
	\end{tabular}
	\label{tab:ablation_study}
\end{table}

\subsection{Ablation Study}
\label{sec:exp_ablation_study}

Our method proposes several steps to handle the uncertainty of sensor data and calibration results, such as the point projection (Section \ref{sec:methodology_fsl_feature_extraction}), uncertainty modelling of normal vectors (Section \ref{sec:methodology_fsc}), and calibration with multiple computations (Section \ref{sec:ext_calibration}).
To show the impact of these individual modules, we conduct the following ablation study.
We disable one of these modules and repeatedly test the proposed method with \textit{RLFS}.
Table \ref{tab:ablation_study} reports the calibration accuracy.

Disabling the point projection \textit{(WO. PP)} or PtL constraints \textit{(WO. PtL)} significantly degrades the calibration performance.
This is because the point projection reduces noise of LiDAR points with the fitted planar parameters and the checkerboard's edges (the source of PtL constraints) offer strong constraints to extrinsics.
We also observe that setting $I=1$ or $I=5$ may sometimes induce unreliable results (see \textit{RLFS03}).
To minimize the error bound of the results, the value of $I$ cannot be small in practice.



\section{Conclusion}\label{sec:conclusion}
In this work, we focus on addressing the fundamental challenge of OOD detection tasks, which is how to fully understand the semantic discrepancy between the ID/OOD samples. We reveal that the key to success in the realistic SCOOD task is to allocate as many ID samples in the unlabeled set correctly as possible. To this end, we propose a novel uncertainty-aware optimal transport scheme that introduces class-specific energy scores as guidance for effective label assignment. Experimental results show that our method achieves better performance than previous state-of-the-art methods on SCOOD benchmarks.

\textbf{Limitations.} In addition to temperature scaling, other techniques such as feature clipping applied in ReAct~\cite{sun2021react} also enhance the performance of energy score, so how to obtain an OOD score that best fits the SCOOD task can be further explored. Moreover, a setting highly related to SCOOD has been proposed in \cite{katz2022training} and formulated as a constrained optimization problem. We will also theoretically analyze these practical OOD settings in our feature work.

% \section*{Acknowledgments}
\textbf{Acknowledgments.} 
This work is supported by National Key R\&D Program of China under Grant 2020AAA0105701, National Natural Science Foundation of China (NSFC) under Grants 61872327, Major Special Science and Technology Project of Anhui, National Natural Science Foundation of China (62033012) and Ant Group through Ant Research Intern Program.


% \section*{Acknowledgement}
% The authors express their sincere gratitude to Prof. Lujia Wang and Mr. Shuyang Zhang for their insightful comments on this article.
% They also wish to extend their heartfelt appreciation to the editorial team of Transactions on Mechatronics and all the reviewers for their prompt and constructive feedback, which significantly enhanced this article.

%%%%%%%%%%%%%%%%
\bibliographystyle{IEEEtran}
\bibliography{reference}{}
% % biography section
% 
% If you have an EPS/PDF photo (graphicx package needed) extra braces are
% needed around the contents of the optional argument to biography to prevent
% the LaTeX parser from getting confused when it sees the complicated
% \includegraphics command within an optional argument. (You could create
% your own custom macro containing the \includegraphics command to make things
% simpler here.)
%\begin{IEEEbiography}[{\includegraphics[width=1in,height=1.25in,clip,keepaspectratio]{mshell}}]{Michael Shell}
% or if you just want to reserve a space for a photo:

\begin{IEEEbiography}[{\includegraphics[width=0.9in,height=1.25in,clip,keepaspectratio]{figure/bio/jjh.png}}]{Jianhao Jiao}
    received the B.Eng. degree in instrument science from Zhejiang University, Hangzhou, China, in 2017, and the Ph.D. from the Department of Electronic and Computer Engineering, the Hong Kong University of Science and Technology, Hong Kong, China, in 2021, supervised by Prof. Ming Liu.
    He is now a research associate in the same university.
    His research interests include state estimation, SLAM, dense mapping, sensor fusion, and computer vision.
\end{IEEEbiography}
\vspace{-0.7cm}

\begin{IEEEbiography}[{\includegraphics[width=0.9in,height=1.15in,clip,keepaspectratio]{figure/bio/cfy_color}}]{Feiyi Chen}
    received the B.Eng. degree from Huazhong University of Science and Technology, Wuhan, China, in 2019 and Master degree from Hong Kong University of Science and Technoloay in 2022, under the supervision of Prof.Ming Liu. His research intersests include multi-sensor fusion, calibration and computer vision.
\end{IEEEbiography}
\vspace{-0.7cm}

\begin{IEEEbiography}[{\includegraphics[width=0.9in,height=1.25in,clip,keepaspectratio]{figure/bio/whx}}]{Hexiang Wei}
    received B.Eng. degree in Mechanical Engineering from Northeastern University, Shenyang,  China, in 2020. He is currently a Ph.D. student in the Intelligent and Autonomous Driving Center at Hong Kong University of Science and Technology led by Prof. Ming Liu. His research interests include SLAM, sensor fusion and robotics.
\end{IEEEbiography}
\vspace{-0.7cm}

\begin{IEEEbiography}[{\includegraphics[width=0.9in,height=1.25in,clip,keepaspectratio]{figure/bio/wujin_color}}]{Jin Wu}
    received the B.S. Degree from University of Electronic Science and Technology of China, Chengdu, China.
    He has been a research assistant since 2018, and now a Ph.D. student in Department of Electronic and Computer Engineering, Hong Kong University of Science and Technology. He was with Tencent Robotics X, Shenzhen, China from 2019 to 2020. During 2012 to
    2018, he was in the UAV industry and has brought up two companies.
\end{IEEEbiography}
\vspace{-0.7cm}

\begin{IEEEbiography}[{\includegraphics[width=0.9in,height=1.25in,clip,keepaspectratio]{figure/bio/ming_color.png}}]{Ming Liu}
    received the Ph.D. degree from the Department of Mechanical and Process Engineering, ETH Zurich, Switzerland, in 2013, supervised by Prof. Roland Siegwart.
    He is currently with the Robotics and Autonomous Systsms, The Hong Kong University of Science and Technology (Guangzhou), the director of Intelligent Autonomous Driving Center, as an Associate Professor.
    He has published many popular papers in top journals including IEEE Transactions on Robotics and International Journal of Robotics Research. Dr. Liu is currently an Associate Editor for IEEE Robotics and Automation Letters, IET Cyber-Systems and Robotics, International Journal of Robotics and Automation.
    His research interests include dynamic environment modeling, deep-learning for robotics, 3D mapping, machine learning, and visual control.
\end{IEEEbiography}



\end{document}




